\section{All orders}
\label{ch:all-orders}

This chapter constructs the cubic action with the whole tower of auxiliary fields $\Sigma_2, \Sigma_3, \ldots$: the action and its equations of motion, the completed squares and the nested kinetic operators, the theorem that fixes the reduced vertices at every order and the collapse lemma behind it, the elimination of the auxiliary fields and the recovered action, the gauge transformations, and the kernel conditions under which the quotient prescription is consistent at every order.
It expands the section of the paper on the all-order construction, which states the same results in condensed form.
The text of that section before its condensation is reproduced in full, followed by the later wording of two of its subsections.
Only this chapter contains the comparison of three bookkeepings of the all-orders problem, the proof of the theorem in the language of the earlier notes, and two complements written from the derivation records: the first orders of the solution of the auxiliary equations, and the converse of the statement that a kernel shift does not change the equation of motion of $\Psi$.
The multiplier form of the action, which the section of the paper also discussed, is in \cref{ch:multiplier-hybrid}, and the reading of the elimination of $\Sigma_2$ as a homotopy transfer is in \cref{ch:ibl}.

\paragraph{Assumptions.}

Everything in this chapter is at tree level and is a statement about formal power series in $\gs$, exact and rational at every truncation of the tower; the convergence of the infinite tower is not addressed.
The algebraic statements hold for symmetric multilinear forms on a finite-dimensional space, which for the string vertices is a level truncation with the momenta fixed or discretized.
Applied to the string vertices, the matching theorem and the results that use it rest on four hypotheses, listed with their status in \cref{ch:hypotheses}.
First, the quartic form $\sprod{2}{2}$ is invertible on the quotient of the level-truncated two-string space by its kernel, and the kernels are stable under the truncation.
Second, the kernel of $\vertex{4}$ vanishes off shell at generic momentum, or, failing that, $\vertex{5}$ and $\vertex{6}$ annihilate it.
Third, there is a presentation of the tower in which every null vector of the kinetic operator of $\Sigma_k$, $k \ge 3$, is a joint null vector of the vertices in which $\Sigma_k$ appears, and in which the reduced sextic vertex satisfies the second kernel clause at $k = 2$.
Fourth, the feedback produced by the exchanges of auxiliary fields is not expressible through the diagonal values $\vertex{n}(\Psi^n)$ and boundary terms; the theorem needs this through the order $\gs^6$.
The first two are supported by the computations at level~0; the last two are open.

\paragraph{Notation across the blocks.}

The blocks taken from the paper before its rewrite write the bare kinetic operator of $\Sigma_k$ as $\kin{k}{}{2 k - 2}$, and the two later blocks write it $\wprod{k}{k}$; the two are the same operator, with $\wprod{2}{2} = \sprod{2}{2}$.
The complements of \cref{ch3:sec:complements} use $\wprod{k}{k}$.

The first block is the section on the all-order construction of the paper as it stood before the condensation, without its last two subsections (the multiplier form and the homotopy transfer, which are in the chapters named above).
Its references to the orders treated one at a time point to \cref{ch:order-by-order}.
It holds the following results.

The rule that $\Sigma_k$ enters at order $\gs^{2 k - 2}$ and its odd coupling at $\gs^{2 k - 1}$ is the definition of the construction, not a result.

\noindent\status{Proven} The minimal count of auxiliary fields, $\Sigma_k$ needed only from $\vertex{3 k - 2}$ on ($k \le \lceil n / 3 \rceil$ for $\vertex{n}$), and the triangular structure of the action: a counting statement, which does not include the dynamical closure.\footnote{Ledger entry R12; its graded-symmetry clause is derived in \nolinkurl{agent/derivations/055-omega-k-graded-symmetry.md}.}

\noindent\status{no ledger entry} The action with the infinite tower as a formal power series, its truncation at a generic order, the equations of motion of $\Psi$ and of every $\Sigma_k$, their order-by-order solution, and the stabilization of the coefficients of the nested operators as the tower grows; the algebra is unconditional, and its application to the string vertices uses the invertibility of the bare operators on their quotients and the kernel condition, and is checked in a finite-dimensional model with the tower $\Sigma_2, \Sigma_3, \Sigma_4$.\record{100-all-orders-action-eom}

\noindent\status{Conditional(PH1, PH2 (b), PH5, PH14 [through $A_6$]) + R70 -- Proven} The matching theorem, the collapse lemma, the value of the action on the solution, the recovered action and the gauge transformations at all orders.
The theorem is proven for generic symmetric tensors in a finite-dimensional space, within the class in which the action is jointly quadratic in the tower and each level enters through one completed square whose kinetic term and source are read off the same form; for the string vertices it holds under the four hypotheses above and the result on the nested kernels stated with the kernel conditions below.
Couplings that break the joint Gaussian form, possible from the octic order on, are not covered, and cubic actions outside the class exist (\cref{ch:multiplier-hybrid,ch:partner-field}).\record{019-b7-allorders-fixed-point (R26)}

\noindent\status{no ledger entry} The equation of motion of the recovered action: the $\Psi$-equation of the cubic action at the auxiliary solution is the equation of motion of the eliminated action, the same for every representative of the solution under the kernel condition; its explicit form is conditional as the matching theorem.\record{085-paper-gaps-3-4}

\subsection{All-order construction}
\label{p:sec:all-orders}



The orders treated one at a time in \cref{p:sec:hs} follow one pattern, and this section states it for every order.
We first say which auxiliary field enters at which order, then write the general form of the cubic action with all the auxiliary fields, its truncation at a generic order and its equations of motion, and explain how it generalizes the squares of \cref{p:sec:hs:quartic,p:sec:hs:sextic,p:sec:hs:octic} and the reduced vertices of \cref{p:sec:hs:sextic,p:sec:hs:septic}.
We then state the theorem that fixes the reduced vertices at every order, and the lemma behind it, that an exchange diagram of auxiliary fields whose sources all lie on the diagonal (i.e.\ the leading order $\Sigma_k \sim \Psi^k$) collapses to a contact term.
The second half of the section eliminates the auxiliary fields one at a time and recovers the original action, gives the gauge transformations and the multiplier form of the action, and reads the elimination of $\Sigma_2$ as a homotopy transfer.
Throughout, we work at tree level and with formal power series in $\gs$: at every truncation of the tower the identities are exact and rational in $\gs$, and we do not address the convergence of the infinite tower.

\subsubsection{Cubic action at all orders}
\label{p:sec:all:action}

\subparagraph{Tower of auxiliary fields.}

The construction of \cref{p:sec:hs} introduces one $k$-string auxiliary field $\Sigma_k \in \Hn{k}$ at each even order: $\Sigma_2$ at order $\gs^2$, $\Sigma_3$ at order $\gs^4$, $\Sigma_4$ at order $\gs^6$.
The rule is that $\Sigma_k$ enters at order $\gs^{2 k - 2}$, through its kinetic term, and that the odd order which follows adds its coupling to $\Psi$:
\begin{equation}
  \label{p:eq:all:counting}
  \Sigma_k
    \;\longleftrightarrow\;
  \Big( \frac{\gs^{2 k - 2}}{(2 k)!} \, \dvertex{2 k}, \; \frac{\gs^{2 k - 1}}{(2 k + 1)!} \, \dvertex{2 k + 1} \Big),
  \qquad
  k \ge 2,
\end{equation}
where $\dvertex{2 k}$ and $\dvertex{2 k + 1}$ are the reduced vertices of these two orders, $\dvertex{4} = \vertex{4}$ and $\dvertex{5} = \vertex{5}$.
The reason is the one given at the septic order.
The only quadratic term available for a $k$-string field is $\dvertex{2 k}(\Sigma_k, \Sigma_k)$, a vertex with $2 k$ strands, and it comes with $\gs^{2 k - 2}$.
Conversely, the even vertex $\vertex{2 k}$ is the first whose square, split $k + k$, has $k$ strands per slot.
Introducing $\Sigma_k$ there is the rule of the construction, one auxiliary field per completed square, and it keeps the action in the triangular form of \cref{p:sec:hs:septic}.

The odd vertex $\dvertex{2 k + 1}$ is then written with the partition $2 k + 1 = k + k + 1$, two slots for $\Sigma_k$ and one for $\Psi$, which keeps the action quadratic in $\Sigma_k$.
It dresses the kinetic operator of $\Sigma_k$ by the background, as the quintic vertex did for $\Sigma_2$ in \cref{p:sec:hs:quintic}.

We call $k$ the level of $\Sigma_k$ in the tower, and keep ``mass level'' for the level of the level truncation.
Each level thus has one pair of consecutive vertices, and the tower of auxiliary fields grows by one field every two orders.
Asking only for a cubic action, not for auxiliary fields eliminated by Gaussian integrals, a cubic action built from the elementary vertices alone could do with fewer fields.
A vertex $\vertex{n}$ can be split as $n = a + b + c$ with $\Sigma_a$, $\Sigma_b$, $\Sigma_c$ on its three legs, so that $\Sigma_k$ is needed only from $\vertex{3 k - 2}$ on, that is $k \le \lceil n / 3 \rceil$ for the vertex $\vertex{n}$.
This count is not the rule of the construction, which introduces $\Sigma_k$ through the completed square~\eqref{p:eq:all:square} at order $\gs^{2 k - 2}$.
The data of level $k$ first enter the value of the action on the solution at the order they are introduced: the even vertex at $\gs^{2 k - 2}$, the odd one at $\gs^{2 k - 1}$, as we show below.

\subparagraph{General form of the action.}

With all the auxiliary fields $\Sigma_2, \Sigma_3, \ldots$, the cubic action is:
  \begin{equation}
  \label{p:eq:all:action}
  \begin{aligned}
    \caction[\infty]
      &
      := S_{\le 3}
        + \sum_{k \ge 2} \Big[
          - \frac{\gs^{2 k - 2}}{(2 k)!} \, \pair[k]{\Sigma_k}{\kin{k}{}{2 k - 2}(\Sigma_k)}
          + \frac{2 \gs^{2 k - 2}}{(2 k)!} \, \pair[k]{\Sigma_k}{\kin{k}{}{2 k - 2}(\Psi \otimes \Sigma_{k - 1})}
      \\ &
      \qquad\qquad\quad
          + \frac{\gs^{2 k - 1}}{(2 k + 1)!} \, \dvertex{2 k + 1}(\Sigma_k, \Sigma_k, \Psi)
        \Big],
    \qquad
    \Sigma_1
      := \Psi,
  \end{aligned}
\end{equation}
where $S_{\le 3}$ is the kinetic term and the cubic vertex of $\Psi$, untouched by the construction, and where the kinetic operator $\kin{k}{}{2 k - 2} \colon \Hn{k} \to \Hn{k}$ of $\Sigma_k$ is read off the reduced even vertex:
\begin{equation}
  \label{p:eq:all:kinetic}
  \begin{gathered}
    \dvertex{2 k}(A, B)
      =: \pair[k]{A}{\kin{k}{}{2 k - 2}(B)},
    \qquad
    A, B \in \Hn{k},
    \\
    \kin{2}{}{2}
      = \sprod{2}{2},
    \qquad
    \kin{3}{}{4}
      = \wprod{3}{3}.
  \end{gathered}
\end{equation}
The pairing is the one of \cref{p:sec:hs:action}, $\pair{A}{B} = \braket{A | c_0^- | B}$,~\eqref{p:eq:hs1:pairing}, and $\pair[k]{\cdot}{\cdot}$ is its $k$-th tensor power, one $c_0^-$ per strand, with the Koszul sign of~\eqref{p:eq:conv:pairing-n}.
For the even tensors met in this section it is the plain product of the pairings of the strands, and it is symmetric.
The auxiliary fields live on the quotients $\Hn{k} / \ker \kin{k}{}{2 k - 2}$.
Each level $k$ of the sum contributes three terms: a kinetic term for $\Sigma_k$, a source for $\Sigma_k$ which is linear in $\Sigma_{k - 1}$ and in $\Psi$, and the odd coupling of $\Sigma_k$ to itself and to $\Psi$.
At $k = 2$ the three terms are those of~\eqref{p:eq:hs4:S4tilde} and~\eqref{p:eq:hs5:S5tilde}, at $k = 3$ those of~\eqref{p:eq:hs7:action-6} and~\eqref{p:eq:hs7:S7tilde}, at $k = 4$ the first two are those of~\eqref{p:eq:hs8:S8tilde} and the third is the coupling $\dvertex{9}(\Sigma_4, \Sigma_4, \Psi)$ named there.
The action is cubic in the fields: apart from $S_{\le 3}$, every term contains two auxiliary fields and at most one $\Psi$, except the quartic source, which contains $\Sigma_2$ and two $\Psi$.

The sum over $k$ is infinite, but~\eqref{p:eq:all:action} is well defined as a formal power series in $\gs$:
  \begin{equation}
  \caction[\infty]
    = \sum_{m \ge 0} \gs^m \, \widetilde{S}_{[m]},
\end{equation}
  whose coefficients are:
  \begin{equation}
  \label{p:eq:all:orders}
  \begin{gathered}
    \widetilde{S}_{[2 k - 2]}
      = \frac{1}{(2 k)!} \Big[
        2 \, \pair[k]{\Sigma_k}{\kin{k}{}{2 k - 2}(\Psi \otimes \Sigma_{k - 1})}
        - \pair[k]{\Sigma_k}{\kin{k}{}{2 k - 2}(\Sigma_k)}
      \Big],
    \\
    \widetilde{S}_{[2 k - 1]}
      = \frac{1}{(2 k + 1)!} \, \dvertex{2 k + 1}(\Sigma_k, \Sigma_k, \Psi),
    \qquad
    k \ge 2,
    \\
    \widetilde{S}_{[0]}
      = \frac{1}{2} \, \pair{\Psi}{Q_B \Psi},
    \qquad
    \widetilde{S}_{[1]}
      = \frac{1}{3!} \, \vertex{3}(\Psi^3).
  \end{gathered}
\end{equation}
  Each order has finitely many terms, two at every even order from $\gs^2$ on and one at every odd order, and each coefficient is a polynomial of degree at most three in finitely many fields.
  The field $\Sigma_k$ enters first at order $\gs^{2 k - 2}$, and the coefficient of $\gs^m$ involves $\Psi$, $\Sigma_{\lfloor m / 2 \rfloor + 1}$ and, for even $m$, $\Sigma_{m / 2}$ only.

We truncate the action at a generic order by truncating this series.
The action with the auxiliary fields $\Sigma_2, \ldots, \Sigma_K$ is the series truncated at order $\gs^{2 K - 1}$, and truncating at the even order $\gs^{2 K - 2}$ instead drops the odd coupling of $\Sigma_K$ alone, as in \cref{p:sec:hs:octic}:
\begin{equation}
\label{p:eq:all:truncation}
\begin{gathered}
  \caction[\le 2 K + 1]
    := \sum_{m = 0}^{2 K - 1} \gs^m \, \widetilde{S}_{[m]},
  \\
  \caction[\le 2 K]
    := \sum_{m = 0}^{2 K - 2} \gs^m \, \widetilde{S}_{[m]}
    = \caction[\le 2 K + 1] - \frac{\gs^{2 K - 1}}{(2 K + 1)!} \, \dvertex{2 K + 1}(\Sigma_K, \Sigma_K, \Psi).
\end{gathered}
\end{equation}
On the solution of the auxiliary equations they reproduce $S_{\le 2 K + 1}$ up to $O(\gs^{2 K})$ and $S_{\le 2 K}$ up to $O(\gs^{2 K - 1})$, the remainder being the order of the first term left out, as we show in~\eqref{p:eq:all:recovered} below.

The action is quadratic in each auxiliary field separately.
The two terms quadratic in $\Sigma_k$ combine into one kinetic operator dressed by the background, the generalization of~\eqref{p:eq:hs5:dressed} and~\eqref{p:eq:hs7:dressed-3}:
\begin{equation}
  \label{p:eq:all:dressed-odd}
  \pair[k]{A}{\kin{k}{\Psi}{2 k - 1}(B)}
    := \pair[k]{A}{\kin{k}{}{2 k - 2}(B)}
      - \frac{\gs}{2 k + 1} \, \dvertex{2 k + 1}(A, B, \Psi),
  \qquad
  A, B \in \Hn{k},
\end{equation}
an operator affine in $\Psi$ and symmetric with respect to the $k$-string pairing.
The action is also quadratic in the tower as a whole, since the sources are bilinear in two consecutive fields.
Collecting the tower in $(\Sigma_2, \ldots, \Sigma_K)$, the truncated action~\eqref{p:eq:all:truncation} reads:
\begin{equation}
  \label{p:eq:all:tridiagonal}
  \caction[\le 2 K + 1]
    = S_{\le 3}
      + \frac{2 \gs^2}{4!} \, \pair[2]{\Sigma_2}{\sprod{2}{2}(\Psi^2)}
      - \sum_{k, l = 2}^{K} \pair[k]{\Sigma_k}{\mathcal K_{k, l; \Psi}(\Sigma_l)},
\end{equation}
with the blocks:
\begin{equation}
  \label{p:eq:all:blocks}
  \begin{gathered}
    \mathcal K_{k, k; \Psi}
      := \frac{\gs^{2 k - 2}}{(2 k)!} \, \kin{k}{\Psi}{2 k - 1},
    \quad
    k = 2, \ldots, K,
    \\
    \mathcal K_{k, k - 1; \Psi}
      := - \frac{\gs^{2 k - 2}}{(2 k)!} \, \kin{k}{}{2 k - 2}(\Psi \otimes \, \cdot \,),
    \quad
    k = 3, \ldots, K,
  \end{gathered}
\end{equation}
with $\mathcal K_{k - 1, k; \Psi}$ the transpose of $\mathcal K_{k, k - 1; \Psi}$ with respect to the pairings,~\eqref{p:eq:conv:transpose}, and $\mathcal K_{k, l; \Psi} = 0$ for $\abs{k - l} \ge 2$.
For the whole series~\eqref{p:eq:all:action} the sum runs over all $k, l \ge 2$, and each coefficient in $\gs$ involves finitely many blocks.
The operator $\mathcal K_\Psi$ is block-tridiagonal: each auxiliary field couples to itself and to its two neighbors only, and the only term linear in the tower is the quartic source of $\Sigma_2$.
Together with sources that depend on the lower-order fields only, this is the triangular form of the action named in \cref{p:sec:hs:septic}, now for the whole tower.

\subparagraph{Equations of motion.}

We normalize the equations of motion as in~\eqref{p:eq:hs8:eom-normalization} of \cref{p:sec:hs:octic}:
\begin{equation}
\delta \caction
  = \pair{\delta \Psi}{\mathcal E_\Psi}
    + \sum_{k} \frac{2 \gs^{2 k - 2}}{(2 k)!} \, \pair[k]{\delta \Sigma_k}{\mathcal E_{\Sigma_k}},
\end{equation}
so that each equation of an auxiliary field has coefficient one on its source, up to the overall normalization of the term in the action.
Varying~\eqref{p:eq:all:action}, or its truncation~\eqref{p:eq:all:truncation}, gives:
\begin{subequations}
\label{p:eq:all:eom}
\begin{align}
  \label{p:eq:all:eom-psi}
  \mathcal E_\Psi
    &
    = Q_B \Psi
      + \frac{\gs}{2} \, \ell_2(\Psi^2)
      + \frac{\gs^2}{3!} \, \ell_3(\Psi, \Sigma_2)
      + \frac{\gs^3}{5!} \, \ell_4(\Sigma_2^2)
  \notag
  \\ &
  \quad
      + \sum_{k = 3}^{K} \Big[
        \frac{2 \gs^{2 k - 2}}{(2 k)!} \, \dvertex{2 k}(\, \cdot \, \otimes \Sigma_{k - 1}, \Sigma_k)
        + \frac{\gs^{2 k - 1}}{(2 k + 1)!} \, \dvertex{2 k + 1}(\Sigma_k, \Sigma_k, \, \cdot \,)
      \Big],
  \\
  \label{p:eq:all:eom-sigma}
  \mathcal E_{\Sigma_k}
    &
    = \kin{k}{}{2 k - 2}(\Psi \otimes \Sigma_{k - 1})
      - \kin{k}{}{2 k - 2}(\Sigma_k)
      + \frac{\gs}{2 k + 1} \, \dvertex{2 k + 1}(\, \cdot \,, \Sigma_k, \Psi)
  \notag
  \\ &
  \quad
      + \frac{\gs^2}{(2 k + 1) (2 k + 2)} \, \dvertex{2 k + 2}(\Psi \otimes \, \cdot \,, \Sigma_{k + 1}),
  \qquad
  2 \le k \le K,
\end{align}
\end{subequations}
for every $K$, the whole series being $K = \infty$, and where the dot denotes the state paired with the open slot, as in~\eqref{p:eq:hs6:eom}.
At $k = 2$ the source is $\sprod{2}{2}(\Psi^2)$ and the odd term is $\tfrac{\gs}{5} \sprod{3}{2}(\Psi, \Sigma_2)$.
The terms of level $2$ in the equation of $\Psi$ are $\ell_3$ and $\ell_4$, the odd term of level $2$.
At $k = K < \infty$ the last term of~\eqref{p:eq:all:eom-sigma} is absent, and in the even truncation the odd term of level $K$ is absent from both equations.
At $K = 2$ this is the $\ell_4$ term, which the quartic equation~\eqref{p:eq:hs4:eom-psi} does not have.
With the dressed operator~\eqref{p:eq:all:dressed-odd}, the equation of $\Sigma_k$ is:
\begin{equation}
\label{p:eq:all:eom-sigma-dressed}
\mathcal E_{\Sigma_k}
  = \kin{k}{}{2 k - 2}(\Psi \otimes \Sigma_{k - 1})
    - \kin{k}{\Psi}{2 k - 1}(\Sigma_k)
    + \frac{\gs^2}{(2 k + 1) (2 k + 2)} \, \dvertex{2 k + 2}(\Psi \otimes \, \cdot \,, \Sigma_{k + 1}),
\end{equation}
the row $k$ of the block-tridiagonal system~\eqref{p:eq:all:tridiagonal}.

The equations of the auxiliary fields are linear in the tower.
The equation of $\Sigma_k$ couples it to the lower-order field $\Sigma_{k - 1}$ through its source, to itself through the operator dressed by the odd vertex, and to the higher-order field $\Sigma_{k + 1}$ through the transpose of the next source, with $\gs^2$.
At finite $K$ the last equation has no higher-order field and closes the chain.
The equation of $\Psi$ is the only non-linear one and the only one that involves every field.
With the tower eliminated, it is the equation of motion of the recovered action,~\eqref{p:eq:all:eom-recovered} below.
Under the condition~\eqref{p:eq:hs8:kernel-condition} the equations are statements on the quotients: $\mathcal E_{\Sigma_k}$ does not depend on the kernel components of the fields and pairs to zero with every kernel direction.
At $K = 4$, with the odd coupling of $\Sigma_4$ dropped, these are the equations~\eqref{p:eq:hs8:eom} of the octic order, written there term by term.

At $K = \infty$, we solve the system order by order in $\gs$.
With
\begin{equation}
  \Sigma_k
    = \sum_m \gs^m \, \Sigma_k^{[m]}
\end{equation}
and $\Sigma_k^{[0]} = \Psi^k$, the coefficient of $\gs^m$ of~\eqref{p:eq:all:eom-sigma} gives:
\begin{equation}
\label{p:eq:all:eom-orders}
\begin{aligned}
  \kin{k}{}{2 k - 2}\big( \Sigma_k^{[m]} \big)
    &
    = \kin{k}{}{2 k - 2}\big( \Psi \otimes \Sigma_{k - 1}^{[m]} \big)
      + \frac{1}{2 k + 1} \, \dvertex{2 k + 1}\big( \, \cdot \,, \Sigma_k^{[m - 1]}, \Psi \big)
  \\ &
  \quad
      + \frac{1}{(2 k + 1) (2 k + 2)} \, \dvertex{2 k + 2}\big( \Psi \otimes \, \cdot \,, \Sigma_{k + 1}^{[m - 2]} \big),
\end{aligned}
\end{equation}
a finite recursion in $(m, k)$, solved on the quotients: each coefficient needs the lower-order field at the same order in $\gs$, the field itself one order down, and the higher-order field two orders down.

\subsubsection{Completed squares}
\label{p:sec:all:squares}

\subparagraph{Nested operators.}

We diagonalize the block-tridiagonal quadratic form by eliminating one field at a time from the top, and each elimination inverts one block.
The operators so obtained are the dressed kinetic operators, defined from the top down, starting from the operator $\kin{K}{\Psi}{2 K - 1}$ of~\eqref{p:eq:all:dressed-odd}, by:
\begin{equation}
  \label{p:eq:all:nested}
  \begin{aligned}
    &
    \pair[k]{A}{\kin{k}{\Psi}{2 K - 1}(B)}
      := \pair[k]{A}{\kin{k}{\Psi}{2 k - 1}(B)}
    \\ &
    \qquad
    - \frac{\gs^2}{(2 k + 1) (2 k + 2)} \, \pair[k + 1]{\kin{k + 1}{}{2 k}(\Psi \otimes A)}{\big( \kin{k + 1}{\Psi}{2 K - 1} \big)^{-1} \circ \kin{k + 1}{}{2 k}(\Psi \otimes B)},
  \end{aligned}
\end{equation}
for $k = K - 1, \ldots, 2$ and $A, B \in \Hn{k}$.
The second term is the exchange of $\Sigma_{k + 1}$ between two sources of $\Sigma_{k + 1}$.
  At $k = 2, 3$, its coefficient is the $1 / 30$ of~\eqref{p:eq:hs6:dressed-2} and the $1 / 56$ of~\eqref{p:eq:hs8:dressed-3}.
The operator of each level contains the inverse of the operator of the next level, and the chain stops at the top, where no higher-order field is left.
We assume, as in \cref{p:sec:hs:octic}, that the bare operators $\kin{k}{}{2 k - 2}$ are invertible on their quotients and that the condition~\eqref{p:eq:hs8:kernel-condition} on the nested kernels holds.
At a fixed level truncation, with the momenta fixed or discretized, and at fixed $K$, the nested operators are then invertible for $\gs \Psi$ in a neighborhood of zero, and they reduce to the bare ones at $\gs \Psi = 0$.
At $K = \infty$ the chain has no top.
Its coefficient of $\gs^m$ involves finitely many levels, as we show below, so that the nested operators are defined coefficient by coefficient, and their inverses exist as formal power series because their term of order $\gs^0$ is the bare operator.
No assumption beyond those of finite $K$ is needed, but we lose the statement that the inverse exists for $\gs \Psi$ in a neighborhood of zero.

The nested operators put the action in a form which is the all-order version of the completed squares of \cref{p:sec:hs}.
Define the fluctuation of each field from its conditional solution, the solution of its own equation of motion at fixed lower-order fields, once the higher-order ones are eliminated:
\begin{equation}
  \label{p:eq:all:deviation}
  \sigma_k^{(2 K - 1)}
    := \Sigma_k - \big( \kin{k}{\Psi}{2 K - 1} \big)^{-1} \circ \kin{k}{}{2 k - 2}(\Psi \otimes \Sigma_{k - 1}),
  \qquad
  k = 2, \ldots, K,
\end{equation}
with $\Sigma_1 = \Psi$, so that $\sigma_2^{(2 K - 1)} = \Sigma_2 - (\kin{2}{\Psi}{2 K - 1})^{-1} \circ \sprod{2}{2}(\Psi^2)$.
Then, exactly and off-shell, we have:
\begin{equation}
  \label{p:eq:all:squares}
  \caction[\le 2 K + 1]
    = \caction[\le 2 K + 1] \big|_{\text{solution}}
      - \sum_{k = 2}^{K} \frac{\gs^{2 k - 2}}{(2 k)!} \, \pair[k]{\sigma_k^{(2 K - 1)}}{\kin{k}{\Psi}{2 K - 1}(\sigma_k^{(2 K - 1)})},
\end{equation}
where the first term is the value of the action on the joint solution of the auxiliary equations, computed below.
At $K = 2$ this is the completed square of \cref{p:sec:hs:quintic},~\eqref{p:eq:hs5:square}, at $K = 3$ the identity~\eqref{p:eq:hs7:identity} of \cref{p:sec:hs:septic}, and the fluctuations generalize $\Sigma_2 - \Psi^2$ at the quartic order and $\Sigma_3 - \Psi \otimes \Sigma_2$ at the sextic one, to which they reduce at $\gs = 0$.
Each square carries the operator of its own level, which is the operator of that field once the higher-order fields are eliminated.
We checked~\eqref{p:eq:all:squares} in the finite-dimensional model, where it holds to machine precision for a tower of three auxiliary fields, $\Sigma_2$, $\Sigma_3$, $\Sigma_4$.

\subparagraph{Squares at every order.}

We reach the form~\eqref{p:eq:all:action} by repeating at every even order the step of \cref{p:sec:hs:sextic,p:sec:hs:octic}.
At order $\gs^{2 k - 2}$ the matching requires a new term, which we write as $\tfrac{\gs^{2 k - 2}}{(2 k)!} \dvertex{2 k}(\Psi \otimes \Sigma_{k - 1}, \Psi \otimes \Sigma_{k - 1})$, with $\dvertex{2 k}$ a symmetric bilinear form on $\Hn{k}$ whose diagonal value the matching fixes.
We complete this term into a square with a new $k$-string field:
\begin{equation}
  \label{p:eq:all:square}
  \begin{gathered}
    \Delta S_{2 k}
      := - \frac{\gs^{2 k - 2}}{(2 k)!} \, \dvertex{2 k}(\Sigma_k - \Psi \otimes \Sigma_{k - 1}, \Sigma_k - \Psi \otimes \Sigma_{k - 1}),
    \\
    \caction[2 k]
      := \frac{\gs^{2 k - 2}}{(2 k)!} \, \dvertex{2 k}(\Psi \otimes \Sigma_{k - 1}, \Psi \otimes \Sigma_{k - 1}) + \Delta S_{2 k}.
  \end{gathered}
\end{equation}
At $k = 2$ this is the square~\eqref{p:eq:hs4:DeltaS4} of $\Sigma_2 - \Psi^2$ with the quartic vertex, at $k = 3$ the square~\eqref{p:eq:hs6:DeltaS6} of $\Sigma_3 - \Psi \otimes \Sigma_2$ with $\dvertex{6}$, at $k = 4$ the square~\eqref{p:eq:hs8:DeltaS8} with $\dvertex{8}$.
Expanding the square cancels the term we started from and leaves the kinetic term and the source of~\eqref{p:eq:all:action}.
The state $\Psi \otimes \Sigma_{k - 1}$, rather than $\Psi^k$, in the square keeps the cross term cubic in the fields.
Since $\Sigma_{k - 1} = \Psi^{k - 1} + O(\gs)$ on-shell, only the diagonal value $\dvertex{2 k}(\Psi^{2 k})$ is produced at this order.
At the odd order which follows, the vertex $\dvertex{2 k + 1}$ is written with two slots for $\Sigma_k$ and one for $\Psi$:
\begin{equation}
  \label{p:eq:all:odd}
  \caction[2 k + 1]
    := \frac{\gs^{2 k - 1}}{(2 k + 1)!} \, \dvertex{2 k + 1}(\Sigma_k, \Sigma_k, \Psi),
\end{equation}
as in~\eqref{p:eq:hs5:S5tilde} and~\eqref{p:eq:hs7:S7tilde}, which is the third term of~\eqref{p:eq:all:action}.

The two terms of level $k$ at the even order come from one form, $\dvertex{2 k}$, which fills the kinetic term with its $2 k$ slots split as $k + k$ and the source with the split $k + (k - 1) + 1$.
This makes $\caction[2 k]$ a square, and the construction rests on this property.
The representative of the reduced vertex built from the diagrams is the elementary vertex minus the exchange diagrams that the lower-order auxiliary fields already produce at that order, with the reduced vertices at their nodes: $\dvertex{6}$ is~\eqref{p:eq:hs6:dressed-vertex}, the sextic vertex minus the $\Sigma_2$ exchange between two quintic vertices, and $\dvertex{7}$ is~\eqref{p:eq:hs7:dressed-vertex}, with a chain of three quintic vertices and a sextic--quintic exchange.
Written with the string vertices at the nodes, the sign of a diagram depends on the diagram: the chain enters $\dvertex{7}$ with the coefficient $+ \tfrac{42}{25}$ of~\eqref{p:eq:hs7:diagonal}, because the diagram $\dvertex{6} \cdot \vertex{5}$, subtracted, contains the chain once more with the opposite sign.
Only its diagonal value is fixed by the matching, through the rule~\eqref{p:eq:hs7:deficit}, and the representative built from the diagrams is one choice among those with that diagonal.

The dressed kinetic operators of~\eqref{p:eq:all:nested} follow the pattern of \cref{p:sec:hs:septic,p:sec:hs:octic} at every level.
Once the fields of higher order than $\Sigma_k$ are eliminated, the operator of $\Sigma_k$ is its bare operator, dressed by the odd vertex at order $\gs$ and by the exchange of $\Sigma_{k + 1}$ at order $\gs^2$.
The $\Sigma_{k + 1}$ line inside that exchange carries the dressed operator of $\Sigma_{k + 1}$ in its turn.
Expanding the inverse in~\eqref{p:eq:all:nested} as a Neumann series, the operator of $\Sigma_2$ is a continued fraction of exchanges, whose depth is the number of auxiliary fields.
The operator $\kin{2}{\Psi}{2 K - 1}$ contains the exchange of $\Sigma_3$, inside which the $\Sigma_3$ line carries the exchange of $\Sigma_4$, and so on up to $\Sigma_K$.
Its expansion in $\gs$ starts as in~\eqref{p:eq:hs8:scaling}, and the coefficient of $\gs^m$ involves the levels $k \le m / 2 + 2$ only, since each rung of the chain costs $\gs^2$.
More precisely, for the action truncated at order $\gs^{2 K - 1}$, the coefficients of $\gs^m$ of the operator of $\Sigma_k$ and of the solution $\Sigma_k^*$ do not depend on $K$ once $K \ge k + \lfloor m / 2 \rfloor$, and that of the value of the action once $K \ge 1 + \lfloor m / 2 \rfloor$.
For the truncation at $\gs^{2 K - 2}$, $\lfloor m / 2 \rfloor$ is replaced by $\lceil m / 2 \rceil$, since the top level has no odd dressing.
These are algebraic identities between the recursions at finite $K$ and at $K = \infty$, valid whenever the bare inverses exist on the quotients.
Every coefficient of the expansion is therefore unchanged once the tower is long enough, and the continued fraction is well defined as a formal power series in $\gs$: the operators at $K = \infty$ are the limits of those at finite $K$, coefficient by coefficient.

The powers of $\gs$ in~\eqref{p:eq:all:action} are those of the original action, in which $\vertex{n}$ carries $\gs^{n - 2}$.
We can redistribute them by absorbing a power of $\gs$ in each auxiliary field, as was done for $\Sigma_2$ in \cref{p:sec:hs:quartic}.
With
\begin{equation}
  \label{p:eq:all:rescaling}
  \widetilde{\Sigma}_k
    := \gs^{k - 1} \, \Sigma_k,
  \qquad
  k \ge 1,
  \qquad
  \widetilde{\Sigma}_1
    = \Psi,
\end{equation}
the action becomes:
\begin{equation}
  \label{p:eq:all:rescaled}
  \begin{aligned}
    \caction[\le 2 K + 1]
      &
      = S_{\le 3}
        + \sum_{k = 2}^{K} \Big[
          - \frac{1}{(2 k)!} \, \pair[k]{\widetilde{\Sigma}_k}{\kin{k}{}{2 k - 2}(\widetilde{\Sigma}_k)}
          + \frac{2 \gs}{(2 k)!} \, \pair[k]{\widetilde{\Sigma}_k}{\kin{k}{}{2 k - 2}(\Psi \otimes \widetilde{\Sigma}_{k - 1})}
      \\ &
      \qquad\qquad\quad
          + \frac{\gs}{(2 k + 1)!} \, \dvertex{2 k + 1}(\widetilde{\Sigma}_k, \widetilde{\Sigma}_k, \Psi)
        \Big].
  \end{aligned}
\end{equation}
Every kinetic term, including the one of $\Psi$ in $S_{\le 3}$, now carries no power of $\gs$, and every cubic term, including the cubic vertex of $\Psi$, carries exactly one, as expected for a cubic action with one coupling constant.
The rescaling reflects a homogeneity of the action: under $\Psi \to t \Psi$, $\Sigma_k \to t^k \Sigma_k$ and $\gs \to \gs / t$, every term of~\eqref{p:eq:all:action} scales as $t^2$, so that the rescaled fields $\widetilde{\Sigma}_k$ all scale as $\Psi$ does.
Since the rescaled fields all have the weight of $\Psi$, their sum is homogeneous under this scaling, and we define the total field:
\begin{equation}
  \label{p:eq:all:total-field}
  \begin{gathered}
    \boldsymbol{\Psi}
      := \sum_{n \ge 1} \gs^{n - 1} \, \Psi_n
      = \sum_{n \ge 1} \widetilde{\Sigma}_n,
    \\
    \Psi_1
      := \Psi,
    \qquad
    \Psi_n
      := \Sigma_n
    \quad
    (n \ge 2),
  \end{gathered}
\end{equation}
an element of $\bigoplus_{n \ge 1} \Hn{n}$.
In terms of it, the kinetic terms of~\eqref{p:eq:all:rescaled} are one quadratic form on the direct sum, with no $\gs$, and the cubic terms are one cubic form with a single $\gs$ in front.
This is the form in which \cref{p:sec:ibl} reads the cubic action as an action for a single string field with multi-string components.

\subsubsection{Reduced vertices at every order}
\label{p:sec:all:reduced-vertices}

The matching of \cref{p:sec:hs} fixed the reduced vertices one order at a time, through their diagonal values.
The following statement extends this to every order.
Let the vertices $\vertex{n}$, $n \ge 4$, be symmetric multilinear forms on a finite-dimensional space, which for the string vertices is a level truncation with the momenta fixed or discretized, and let every identity be read as a formal power series in $\gs$.
Assume that the bare kinetic operators $\kin{k}{}{2 k - 2}$ are invertible on their quotients and that the condition~\eqref{p:eq:hs8:kernel-condition} on the nested kernels holds.
The nested operators~\eqref{p:eq:all:nested} are then invertible for $\gs \Psi$ in a neighborhood of zero at every finite $K$, the neighborhood possibly depending on $K$, and as formal power series at $K = \infty$.
Assume also that each level is introduced through the square~\eqref{p:eq:all:square}, its kinetic term and its source read off the same form $\dvertex{2 k}$.
Then there exist reduced vertices $\dvertex{n}$, $n \ge 4$, with $\dvertex{4} = \vertex{4}$ and $\dvertex{5} = \vertex{5}$, whose diagonal values are fixed recursively by:
\begin{equation}
  \label{p:eq:all:diagonal}
  \dvertex{n}(\Psi^n)
    = \vertex{n}(\Psi^n) - n! \, \big[ \gs^{n - 2} \big] \, \caction[\le n - 1] \big|_{\text{solution}},
  \qquad
  n \ge 6,
\end{equation}
such that the cubic action reproduces the original one on the solution of the auxiliary equations:
\begin{equation}
  \label{p:eq:all:closure}
  \caction[\le N] \big|_{\text{solution}}
    = S_{\le N} + O(\gs^{N - 1}),
  \qquad
  N \ge 4,
\end{equation}
as a formal power series in $\gs$.
At each truncation the identity is exact and rational in $\gs$, with a remainder at the order of the first term left out, $\gs^{2 K}$ for $\caction[\le 2 K + 1]$ and $\gs^{2 K - 1}$ for $\caction[\le 2 K]$.
The rule~\eqref{p:eq:all:diagonal} is~\eqref{p:eq:hs7:deficit}: the diagonal value of the new vertex is the vertex of the original action minus what the previous orders already produce, and~\eqref{p:eq:hs6:dressed-vertex}, \eqref{p:eq:hs7:coefficient} and~\eqref{p:eq:hs8:diagonal} are its cases $n = 6, 7, 8$.

Three facts make the recursion well defined.
First, the mismatch on the right of~\eqref{p:eq:all:diagonal} is a homogeneous polynomial of degree $n$ in $\Psi$: a polynomial because finitely many diagrams contribute to it, and homogeneous by the homogeneity under $\Psi \to t \Psi$, $\Sigma_k \to t^k \Sigma_k$, $\gs \to \gs / t$ noted above, since the coefficient of $\gs^{n - 2}$ in a quantity scaling as $t^2$ has degree $n$.
A homogeneous polynomial of degree $n$ is the diagonal value of a symmetric $n$-linear form, its polarization, which we read as a form on $\Hn{k} \times \Hn{k}$ for $n = 2 k$ and on $\Hn{k} \times \Hn{k} \times \Hs$ for $n = 2 k + 1$ by grouping its slots, with the same value at $\Psi^n$; a vertex with the required diagonal therefore exists.
Second, the coefficient of $\gs^{n - 2}$ in the eliminated action depends on $\dvertex{n}$ only through its diagonal value, with coefficient $1 / n!$, and on the vertices $\dvertex{m}$ with $m < n$ otherwise.
This is the content of the collapse lemma below.
The diagonal of $\dvertex{n}$ is therefore fixed by the orders below it, and fixing it does not disturb them.

Third, everything else in $\dvertex{n}$ is free as far as the matching goes: two forms with the same diagonal value close the order equally well.
Where the kinetic operators have a kernel, the condition~\eqref{p:eq:hs8:kernel-condition} restricts this freedom (\cref{p:sec:hs:octic}).
This freedom has an effect on the next orders: the terms of $\dvertex{n}$ vanishing on the diagonal enter their exchange diagrams and change the mismatches there, hence the diagonal values of the later vertices.
The recursion absorbs the change by fixing those diagonals in their turn.
At the sextic order, for instance, the vertex~\eqref{p:eq:hs6:dressed-vertex} built from the exchange diagram and the polarization of its diagonal value both close the order and the next ones, differ as forms, and give different diagonal values for $\dvertex{7}$.

We prove the statement by induction on $n$ at the algebraic level, for symmetric forms on a finite-dimensional space.
The induction hypothesis is that the vertices $\dvertex{m}$, $m < n$, exist and close the orders below $\gs^{n - 2}$; the first fact gives a $\dvertex{n}$ with the required diagonal, and the second, which rests on the collapse lemma below, shows that it closes the order $\gs^{n - 2}$ without changing the lower ones.
We ran the recursion in the finite-dimensional model from $n = 6$ to $n = 9$: the remainder of the eliminated action scales as $\gs^8$, and its response to a change of the diagonal is the one stated below.
No property of the vertices beyond their symmetry is used: the relations satisfied by the string vertices play no role in the matching, which is a statement about the Hubbard--Stratonovich transformation alone.
From the string vertices we need in addition the invertibility assumption on the kinetic operators, which we take as in \cref{p:sec:hs:octic}.
The independence of the exchange structures is needed only to decompose the diagonal into named terms, as in~\eqref{p:eq:hs7:diagonal}.
The recursion itself fixes the diagonal as a polynomial and needs no such independence.

The statement is one about formal power series: at every truncation the identities are exact, and nothing is said about the convergence of the continued fraction~\eqref{p:eq:all:nested} as the number of auxiliary fields grows.
We do not write a representative of $\dvertex{n}$ built from the diagrams for $n \ge 8$ either, as noted in \cref{p:sec:hs:octic}.

\subparagraph{Value of the action on the solution.}

With the nested operators built from the top down, the auxiliary equations of~\eqref{p:eq:all:truncation} are solved from the bottom up, each field in terms of the lower-order ones:
\begin{equation}
  \label{p:eq:all:saddles}
  \Sigma_k^*
    = \big( \kin{k}{\Psi}{2 K - 1} \big)^{-1} \circ \kin{k}{}{2 k - 2}(\Psi \otimes \Sigma_{k - 1}^*),
  \qquad
  k = 2, \ldots, K,
  \qquad
  \Sigma_1^*
    := \Psi,
\end{equation}
the fluctuations~\eqref{p:eq:all:deviation} vanishing one after the other.
Each solution lies on the diagonal at leading order, $\Sigma_k^* = \Psi^k + O(\gs)$, as we show below.
The squares of~\eqref{p:eq:all:squares} vanish on the joint solution, and the value of the action is the first term there.
With the auxiliary fields eliminated one after the other, in decreasing order, $\Sigma_K$ first, this value is:
\begin{equation}
  \label{p:eq:all:value}
  \caction[\le 2 K + 1] \big|_{\text{solution}}
    = S_{\le 3} + \frac{\gs^2}{4!} \, \pair[2]{\sprod{2}{2}(\Psi^2)}{\big( \kin{2}{\Psi}{2 K - 1} \big)^{-1} \circ \sprod{2}{2}(\Psi^2)},
\end{equation}
exactly in $\gs$ at fixed truncation, under the invertibility assumption stated above.
The only source is the quartic one, and every other term of the action sits in the nested operator~\eqref{p:eq:all:nested} of $\Sigma_2$, a continued fraction in which the dressed operator of each field contains the inverse of the next one.
At $K = 2$ the operator is $\kin{2}{\Psi}{3}$ and~\eqref{p:eq:all:value} is the constant term of~\eqref{p:eq:hs5:square} in \cref{p:sec:hs:quintic}; with the three auxiliary fields of \cref{p:sec:hs:octic} and the top odd coupling dropped, it is the value of that section with the operator~\eqref{p:eq:hs8:dressed-2}.
As a formal series in $\gs$, every coefficient of the expansion of~\eqref{p:eq:all:value} no longer changes once enough fields are introduced, and the whole interaction content of the original action is carried by the diagonal values of the reduced vertices.
At $K = \infty$, under the assumptions of the matching statement above, the value of the action~\eqref{p:eq:all:action} on the solution is then the original action $S$, as a formal power series in $\gs$.
We write out the elimination itself, with the induced terms at each step, below.

The lemma behind~\eqref{p:eq:all:diagonal} is the all-order form of the collapse of \cref{p:sec:hs:quartic}, where the $\Sigma_2$ exchange between two quartic sources gave back the contact term $\vertex{4}(\Psi^4)$,~\eqref{p:eq:hs4:telescoping}.
It has three parts.
First, because the kinetic term and the source of level $k$ come from the same form, the source evaluated on the diagonal is the kinetic operator applied to the diagonal: for $A \in \Hn{k}$, we have:
\begin{equation}
  \label{p:eq:all:same-form}
  \pair[k]{A}{\kin{k}{}{2 k - 2}(\Psi \otimes \Psi^{k - 1})}
    = \dvertex{2 k}(A, \Psi^k)
    = \pair[k]{A}{\kin{k}{}{2 k - 2}(\Psi^k)},
\end{equation}
where the first reading is the source of~\eqref{p:eq:all:action} at $\Sigma_{k - 1} = \Psi^{k - 1}$ and the last one its kinetic term at $\Sigma_k = \Psi^k$, both being $\dvertex{2 k}$ with all slots but the first $k$ filled with $\Psi$.
Hence $(\kin{k}{}{2 k - 2})^{-1}$ applied to the source on the diagonal gives $\Psi^k$ on the quotient, and since the dressed operators reduce to the bare ones at $\gs = 0$, the solutions~\eqref{p:eq:all:saddles} are $\Sigma_k^* = \Psi^k + O(\gs)$ by induction on $k$.
No symmetry of $\dvertex{2 k}$ beyond the one in the two groups of $k$ slots is needed, so that the forms built from the exchange diagrams and the polarized ones qualify equally.
Second, each rung of the chain~\eqref{p:eq:all:nested} costs $\gs^2$, so that the data of level $k$ first enter the value~\eqref{p:eq:all:value} at order $\gs^2 (\gs^2)^{k - 2} = \gs^{2 k - 2}$.
The odd vertex $\dvertex{2 k + 1}$, which dresses the operator of level $k$ at order $\gs$, enters one order later.
Third, at its first order each level contributes its diagonal value alone, with the coefficient of the original action: at that order the two sources flanking the level-$k$ exchange are on the diagonal, and
\begin{equation}
  \label{p:eq:all:collapse}
  \frac{\gs^{2 k - 2}}{(2 k)!} \, \pair[k]{\kin{k}{}{2 k - 2}(\Psi^k)}{\big( \kin{k}{}{2 k - 2} \big)^{-1} \circ \kin{k}{}{2 k - 2}(\Psi^k)}
    = \frac{\gs^{2 k - 2}}{(2 k)!} \, \dvertex{2 k}(\Psi^{2 k}),
\end{equation}
the inverse cancelling against one of the two operators, as $\sprod{2}{2}^{-1}$ did in~\eqref{p:eq:hs4:telescoping}.
The coefficient $1 / (2 k)!$ is the product of the $1 / 4!$ of the quartic source and of the factors $(2 j)! / (2 j + 2)!$ of the rungs $j = 2, \ldots, k - 1$ in~\eqref{p:eq:all:nested}, which telescope.
For the odd vertex, expanding $(\kin{k}{\Psi}{2 k - 1})^{-1}$ to first order in $\gs$ between the same two diagonal sources gives $+ \tfrac{\gs^{2 k - 1}}{(2 k + 1)!} \dvertex{2 k + 1}(\Psi^{2 k + 1})$, as $\vertex{5}$ entered in~\eqref{p:eq:hs5:constant}.
An exchange diagram of auxiliary fields with all its sources on the diagonal thus collapses to a contact term, at every level.
From the sextic order on, the exchanges between interaction vertices do not collapse, and the reduced vertices absorb them.
In the finite-dimensional model, adding to $\dvertex{8}$ the form $\epsilon \, \alpha^{\otimes 8}$, with $\alpha$ a linear form, shifts the coefficient of $\gs^6$ in the value by $\epsilon \, \alpha(\Psi)^8 / 8!$ and leaves the lower orders unchanged, and likewise for $\dvertex{9}$ at order $\gs^7$.

\subsubsection{Integrating out the auxiliary fields}
\label{p:sec:all:integrating-out}

We now carry out the elimination which the value~\eqref{p:eq:all:value} summarizes, one field at a time, to show how the original action comes back and where each induced interaction goes.
The order is from the top of the tower down, since $\Sigma_K$ is the only field whose equation involves no higher-order one.
Each step is exact, because the action is quadratic in the field being eliminated: the elimination of $\Sigma_k$ is the classical form of a Gaussian integral, and it produces no remainder of higher degree in the remaining fields.
A coupling cubic in an auxiliary field would break this: its square never completes, and the elimination would leave a remainder cubic in the fluctuation.
The form with $\Sigma_2$ alone, mentioned in \cref{p:sec:hs:sextic}, has such a coupling, its sextic vertex with three $\Sigma_2$ slots.

The first step eliminates $\Sigma_K$.
In the truncated action~\eqref{p:eq:all:truncation} it appears in the three terms of level $K$ only, and its equation of motion is the case $k = K$ of~\eqref{p:eq:all:eom-sigma-dressed}, which has no higher-order field:
\begin{equation}
\label{p:eq:all:eom-top}
\mathcal E_{\Sigma_K}
  = \kin{K}{}{2 K - 2}(\Psi \otimes \Sigma_{K - 1}) - \kin{K}{\Psi}{2 K - 1}(\Sigma_K)
  = 0,
\end{equation}
the source minus the dressed kinetic operator~\eqref{p:eq:all:dressed-odd} applied to the field, as~\eqref{p:eq:hs5:eom-sigma-dressed} at the quintic order.
Its solution is the top line of~\eqref{p:eq:all:saddles}, $\Sigma_K^* = (\kin{K}{\Psi}{2 K - 1})^{-1} \circ \kin{K}{}{2 K - 2}(\Psi \otimes \Sigma_{K - 1})$, a function of $\Sigma_{K - 1}$ and $\Psi$.
Substituting it back, we find that the kinetic term and the source of level $K$ combine into one term:
\begin{equation}
  \label{p:eq:all:eliminate-top}
  \begin{aligned}
    &
    \caction[\le 2 K + 1] \big|_{\Sigma_K = \Sigma_K^*}
      = \caction[\le 2 K - 1]
    \\ &
    \qquad
    + \frac{\gs^{2 K - 2}}{(2 K)!} \, \pair[K]{\kin{K}{}{2 K - 2}(\Psi \otimes \Sigma_{K - 1})}{\big( \kin{K}{\Psi}{2 K - 1} \big)^{-1} \circ \kin{K}{}{2 K - 2}(\Psi \otimes \Sigma_{K - 1})},
  \end{aligned}
\end{equation}
where $\caction[\le 2 K - 1]$ is the action of the levels $2, \ldots, K - 1$.
The second term is the interaction induced by the $\Sigma_K$ exchange: two sources of $\Sigma_K$ joined by the dressed $\Sigma_K$ line.
It is exactly quadratic in $\Sigma_{K - 1}$, since the source is linear in it, and it carries one $\Psi$ on each side.
At $K = 3$ this is the elimination of $\Sigma_3$ in \cref{p:sec:hs:septic}, whose induced interaction is quadratic in $\Sigma_2$ and combines with the quintic coupling into a single quadratic form.

The induced term has the form of the quadratic term of $\Sigma_{K - 1}$, and we absorb it there.
The sum of the two is:
\begin{equation}
  \label{p:eq:all:absorb}
  \begin{aligned}
    &
    - \frac{\gs^{2 K - 4}}{(2 K - 2)!} \, \pair[K - 1]{\Sigma_{K - 1}}{\kin{K - 1}{\Psi}{2 K - 3}(\Sigma_{K - 1})}
      + \big( \text{induced term of~\eqref{p:eq:all:eliminate-top}} \big)
    \\ &
    \qquad
    = - \frac{\gs^{2 K - 4}}{(2 K - 2)!} \, \pair[K - 1]{\Sigma_{K - 1}}{\kin{K - 1}{\Psi}{2 K - 1}(\Sigma_{K - 1})},
  \end{aligned}
\end{equation}
with the nested operator of~\eqref{p:eq:all:nested}.
At this point the action of the levels $2, \ldots, K - 1$ has the form~\eqref{p:eq:all:truncation} again, with one field less and with the kinetic operator of the new top field dressed by the exchange of the field just eliminated.
The second step eliminates $\Sigma_{K - 1}$ with this dressed operator, by the same two displays with $K$ replaced by $K - 1$ and $\kin{K - 1}{\Psi}{2 K - 3}$ replaced by $\kin{K - 1}{\Psi}{2 K - 1}$, and so on down to $\Sigma_2$.
The chain of operators~\eqref{p:eq:all:nested} is built in the process, one rung per step, and the last step leaves the value~\eqref{p:eq:all:value}: the quartic source of $\Sigma_2$ paired with itself through the fully dressed operator $\kin{2}{\Psi}{2 K - 1}$.
Every step inverts an operator of the form bare plus $O(\gs \Psi)$, so the result is a rational function of $\gs$ at a fixed truncation of the tower and of the mass level, with the momenta fixed or discretized.
The whole tower of induced interactions is this one function.
We checked in the finite-dimensional model, for a tower of three auxiliary fields, that the elimination one field at a time agrees with the simultaneous solution of the auxiliary equations, at real and complex values of $\gs$, to machine precision.

The original action is recovered from the value by the collapse of the exchange diagrams.
If we expand~\eqref{p:eq:all:value} in $\gs$, every term is a diagram of $\Psi$ lines and auxiliary lines with the sources and the odd vertices at its nodes.
The diagrams whose sources all sit on the diagonal collapse to contact terms by~\eqref{p:eq:all:collapse}.
At order $\gs^{2 k - 2}$ the $\Sigma_k$ exchange between two diagonal sources gives $\tfrac{\gs^{2 k - 2}}{(2 k)!} \dvertex{2 k}(\Psi^{2 k})$, and at order $\gs^{2 k - 1}$ the odd vertex inserted on the $\Sigma_k$ line gives $\tfrac{\gs^{2 k - 1}}{(2 k + 1)!} \dvertex{2 k + 1}(\Psi^{2 k + 1})$.
The remaining diagrams at order $\gs^{n - 2}$ are the exchanges between interaction vertices, those which \cref{p:sec:hs:sextic,p:sec:hs:septic} met at the sextic and septic orders.
These are exactly the terms by which $\dvertex{n}(\Psi^n)$ differs from $\vertex{n}(\Psi^n)$ in~\eqref{p:eq:all:diagonal}.
The diagonal value of the reduced vertex was defined as that of the elementary one minus what the lower-order fields already produce, so that at every order the contact term and the exchange terms add up to $\tfrac{\gs^{n - 2}}{n!} \vertex{n}(\Psi^n)$.
Order by order the induced interactions are therefore the terms that the reductions subtract, and the value is:
\begin{equation}
  \label{p:eq:all:recovered}
  \caction[\le 2 K + 1] \big|_{\text{solution}}
    = S_{\le 2 K + 1} + O(\gs^{2 K}),
\end{equation}
the closure~\eqref{p:eq:all:closure} at $N = 2 K + 1$, with the remainder at the order of the first level left out.
The Hubbard--Stratonovich construction thus gives back the original action not term by term, but through the two mechanisms of \cref{p:sec:hs}, repeated at every level: the collapse of the diagonal exchanges and the subtraction of the others.

The equations of motion follow the same pattern.
Varying the cubic action with respect to $\Psi$ and evaluating at the auxiliary solution gives the equation of motion of the value~\eqref{p:eq:all:recovered}, since the auxiliary fields are stationary there and their variation drops out.
Hence, as a formal power series in $\gs$, the $\Psi$ equation of the cubic action with the auxiliaries eliminated is the equation of motion of $S_{\le 2 K + 1}$ up to order $\gs^{2 K - 1}$:
\begin{equation}
  \label{p:eq:all:eom-recovered}
  \mathcal E_\Psi \big|_{\Sigma_k = \Sigma_k^*}
    = Q_B \Psi + \sum_{n = 2}^{2 K} \frac{\gs^{n - 1}}{n!} \, \ell_n(\Psi^n) + O(\gs^{2 K}),
\end{equation}
with the normalization of the products fixed in \cref{p:sec:hs:action}.
The auxiliary solutions~\eqref{p:eq:all:saddles} are determined only modulo the kernels of the bare kinetic operators $\kin{k}{}{2 k - 2}$, and the statement does not depend on the representative.
The assumptions are those of the matching statement above: the bare operators invertible on their quotients, the condition on the nested kernels, and the kinetic term and the source of each level read off one form.
The statement is made at a fixed level truncation with the momenta fixed or discretized.
This is the form in which \cref{p:sec:bar-cobar} uses the result.

Under the condition~\eqref{p:eq:hs8:kernel-condition} the auxiliary solution is stationary in every direction, those of the kernels included, so the variation of $\Sigma_k^*$ drops out whatever its direction.
The remainder of~\eqref{p:eq:all:recovered} is of order $\gs^{2 K}$ and of the form $\gs^{-2}$ times a function of $\gs \Psi$, so its $\Psi$-derivative is of the same order, which is the error term of~\eqref{p:eq:all:eom-recovered}.
A shift $\Sigma_k \to \Sigma_k + N_k$ with $N_k \in \ker \kin{k}{}{2 k - 2}$ leaves the kinetic term of $\Sigma_k$ and its own source unchanged, by the symmetry of $\kin{k}{}{2 k - 2}$, and changes the action by the shifts of the odd coupling and of the source of $\Sigma_{k + 1}$:
\begin{equation}
  \label{p:eq:all:kernel-shift}
  \begin{aligned}
    &
    \frac{\gs^{2 k - 1}}{(2 k + 1)!} \Big[
      2 \dvertex{2 k + 1}(N_k, \Sigma_k, \Psi)
      + \dvertex{2 k + 1}(N_k, N_k, \Psi)
    \Big]
    \\ &
    \qquad
    + \frac{2 \gs^{2 k}}{(2 k + 2)!} \, \pair[k + 1]{\Sigma_{k + 1}}{\kin{k + 1}{}{2 k}(\Psi \otimes N_k)},
  \end{aligned}
\end{equation}
the last term being absent at $k = K$.
These are the only terms through which a kernel shift reaches $\mathcal E_\Psi$, and they are linear in $\Psi$.
The condition~\eqref{p:eq:hs8:kernel-condition} makes each of them vanish, so the $\Psi$-equation, the auxiliary equations and the value are the same for every representative.

\subsubsection{Gauge transformations at all orders}
\label{p:sec:all:gauge}

The gauge transformation~\eqref{p:eq:hs1:gauge-psi} of the string field extends to the tower by the rule of \cref{p:sec:hs:quintic,p:sec:hs:septic}: each auxiliary field transforms as its solution, plus a rotation of its fluctuation which keeps the corresponding square invariant.
Writing the action as a sum of squares,~\eqref{p:eq:all:squares}, makes the rule explicit at every level at once.
Its first term is a function of $\Psi$ alone, and each square is a quadratic form in the fluctuation $\sigma_k^{(2 K - 1)}$ with the field-dependent operator $\kin{k}{\Psi}{2 K - 1}$.
Demanding that every square be exactly invariant under an arbitrary variation $\delta_\Lambda \Psi$ gives, as in~\eqref{p:eq:hs7:gauge-deviations}:
\begin{equation}
  \label{p:eq:all:gauge-deviations}
  \delta_\Lambda \sigma_k^{(2 K - 1)}
    = - \frac{1}{2} \, \big( \kin{k}{\Psi}{2 K - 1} \big)^{-1} \circ \big( \delta_\Lambda \kin{k}{\Psi}{2 K - 1} \big) (\sigma_k^{(2 K - 1)}),
  \qquad
  k = 2, \ldots, K,
\end{equation}
and the field itself transforms as:
\begin{equation}
  \label{p:eq:all:gauge-fields}
  \delta_\Lambda \Sigma_k
    = \delta_\Lambda \Big[ \big( \kin{k}{\Psi}{2 K - 1} \big)^{-1} \circ \kin{k}{}{2 k - 2}(\Psi \otimes \Sigma_{k - 1}) \Big]
      + \delta_\Lambda \sigma_k^{(2 K - 1)},
\end{equation}
the first term being the variation of its conditional solution~\eqref{p:eq:all:deviation}, through $\Psi$ and through $\Sigma_{k - 1}$, which transforms in its turn.
The variation of the nested operator is read off~\eqref{p:eq:all:dressed-odd} and~\eqref{p:eq:all:nested}.
It collects the variation of the odd coupling, $- \tfrac{\gs}{2 k + 1} \dvertex{2 k + 1}(\, \cdot \,, \, \cdot \,, \delta_\Lambda \Psi)$, those of the two sources of $\Sigma_{k + 1}$ flanking the exchange, and that of the inverse of $\kin{k + 1}{\Psi}{2 K - 1}$.
The last one contains the same three pieces one level up.
The laws are rotations rather than implicit equations, since each square is Gaussian in its fluctuation.
With them, the variation of the cubic action equals the variation of its value on the solution identically, off-shell, for any $\delta_\Lambda \Psi$ and exactly in $\gs$ at fixed truncation:
\begin{equation}
  \label{p:eq:all:delta-S-equal}
  \delta_\Lambda \caction[\le 2 K + 1]
    = \delta_\Lambda \Big[ \caction[\le 2 K + 1] \big|_{\text{solution}} \Big],
\end{equation}
so that by~\eqref{p:eq:all:recovered} the cubic action is invariant up to order $\gs^{2 K}$, the order of the truncation it reproduces.
This last step uses the invariance of $S_{\le 2 K + 1}$ to that order, which rests on the relations of \cref{p:sec:hs:action} among the string vertices, whereas the matching itself does not use them.
At $K = 2$ these are the laws~\eqref{p:eq:hs5:gauge-deviation} and~\eqref{p:eq:hs5:gauge-sigma}, at $K = 3$ the laws~\eqref{p:eq:hs7:gauge-deviations} and~\eqref{p:eq:hs7:gauge-fields}.
We checked the identity~\eqref{p:eq:all:delta-S-equal} in the finite-dimensional model for the tower $\Sigma_2, \Sigma_3, \Sigma_4$, with analytic directional derivatives.
Dropping the rotation of $\sigma_2^{(2 K - 1)}$ or of $\sigma_4^{(2 K - 1)}$ leaves exactly the variation of the corresponding square, as it should.

Two features of~\eqref{p:eq:all:gauge-deviations} carry over from the lower orders.
The rotation is homogeneous in $\sigma_k^{(2 K - 1)}$, and the laws carry no gauge parameter of their own at any level: the auxiliary sector follows the solution and rotates its fluctuations.
In the Batalin--Vilkovisky language of \cref{p:sec:gauge}, each fluctuation and its antifield form a contractible pair whose coefficient depends on $\Psi$.
A flat representative makes the coefficient constant.
We constructed it only through order $\gs^4$, but we expect it to exist at any order.
Below the top level the dependence comes through the odd coupling of the level and through the exchanges of the higher-order fields.
At the top level, it comes through the odd coupling alone.
The rotation vanishes when the operator of the level does not depend on $\Psi$, as for $\Sigma_3$ at the sextic order,~\eqref{p:eq:hs6:gauge-sigma3}.

The closure of the gauge algebra, and the representative of the law in which it closes modulo the equation of motion of $\Psi$ alone at the quintic order, are the subject of \cref{p:sec:gauge}.
So is the question whether the cubic action is the gauge-fixed form of a theory with an independent $2$-string parameter: within the classes examined there, every such theory is a trivial enlargement, in which the parameter shifts $\Sigma_2$ along null directions and carries no physics, and the question is left open beyond them.



\medskip
The next two blocks are the subsections on the recovered action and on the kernel conditions as the paper states them after its rewrite, which moved the kernel conditions from the octic order to this section and wrote them for every level.
For these two statements the later wording is the authoritative one.
The earlier wording of the recovered action, in the block above, is kept because it carries the elimination of the fields one at a time and the chain-rule argument for the equation of motion step by step.
The first of the two blocks restates the recovered action~\eqref{p:eq:all:recovered}, its equation of motion and the kernel shift.

\subsubsection{Recovered action and equations of motion}
\label{pc:sec:all:integrating-out}

Expanding the value~\eqref{p:eq:all:value} in $\gs$ shows how the original action comes back.
Every term of the expansion is a tree diagram with $\Psi$ lines and auxiliary lines, the sources and the odd vertices at its nodes.
The diagrams whose sources all sit on the diagonal collapse to contact terms by the lemma: at order $\gs^{2 k - 2}$ the $\Sigma_k$ exchange between two diagonal sources gives $\dvertex{2 k}(\Psi^{2 k})$, and at order $\gs^{2 k - 1}$ the odd vertex on the $\Sigma_k$ line gives $\dvertex{2 k + 1}(\Psi^{2 k + 1})$, both with the coefficients of the original action.
The remaining diagrams at order $\gs^{n - 2}$ are the exchanges between interaction vertices, and by~\eqref{p:eq:all:diagonal} they are the terms by which $\dvertex{n}(\Psi^n)$ differs from $\vertex{n}(\Psi^n)$.
The contact term and the exchange terms therefore add up to $\vertex{n}(\Psi^n)$ at every order:
\begin{equation}
  \label{pc:eq:all:recovered}
  \caction[\le 2 K + 1] \big|_{\Sigma = \Sigma^*}
    = S_{\le 2 K + 1} + O(\gs^{2 K}),
\end{equation}
which is the closure~\eqref{p:eq:all:closure} at $N = 2 K + 1$, and for the full tower the remainder is absent.
The construction gives back the original action through the two mechanisms of \cref{p:sec:hs} repeated at every level: the collapse of the diagonal exchanges and the subtraction of the others.

The equation of motion of $\Psi$ follows.
Since the auxiliary fields are stationary on their solution, their variation drops out, and the $\Psi$ equation of the cubic action with the tower eliminated is the equation of motion of the value~\eqref{pc:eq:all:recovered}:
\begin{equation}
  \label{pc:eq:all:eom-recovered}
  \mathcal E_\Psi \big|_{\Sigma = \Sigma^*}
    = Q_B \Psi + \sum_{n = 2}^{2 K} \frac{\gs^{n - 1}}{n!} \, \ell_n(\Psi^n) + O(\gs^{2 K}),
\end{equation}
with the products normalized as in \cref{p:sec:hs:action}; the remainder of the value is $\gs^{-2}$ times a function of $\gs \Psi$, so that its derivative has the same order.
\Cref{p:sec:bar-cobar} uses the result in this form.

The solutions~\eqref{p:eq:all:saddles} are determined only modulo the kernels of the bare operators, and the statements above do not depend on the representative.
A shift $\Sigma_k \to \Sigma_k + N_k$ with $N_k \in \ker \wprod{k}{k}$ leaves the kinetic term of $\Sigma_k$ and its own source unchanged, by the symmetry of $\wprod{k}{k}$, and changes the action by the shifts of the odd coupling and of the source of $\Sigma_{k + 1}$:
\begin{equation}
  \label{pc:eq:all:kernel-shift}
  \begin{aligned}
    &
    \frac{\gs^{2 k - 1}}{(2 k + 1)!} \Big[
      2 \, \dvertex{2 k + 1}(N_k, \Sigma_k, \Psi)
      + \dvertex{2 k + 1}(N_k, N_k, \Psi)
    \Big]
    \\ &
    \qquad
    + \frac{2 \gs^{2 k}}{(2 k + 2)!} \, \pair[k + 1]{\Sigma_{k + 1}}{\wprod{k + 1}{k + 1}(\Psi \otimes N_k)},
  \end{aligned}
\end{equation}
the last term being absent at the top of a truncated tower.
These are the only terms through which a kernel shift reaches the action and its equations.
We assume that the kernel conditions of \cref{pc:sec:all:kernel} hold, so that these terms vanish: the equations and the value are the same for every representative, and the solution is stationary in every direction, the kernel directions included.




\medskip
The second block states the kernel conditions for every level.
Its display is the condition~\eqref{p:eq:hs8:kernel-condition} of the octic order, written with the later notation; its first clause at $k = 2$ is the condition of the quintic order, and its second clause at $k = 2$ the condition~\eqref{p:eq:hs6:kernel-condition} of the sextic order.
The footnote at its end summarizes the multiplier and hybrid forms, which are treated in full in \cref{ch:multiplier-hybrid}.

\noindent\status{mixed} The quotient of the two-string space by the kernel of $\vertex{4}$ is compatible with the quintic vertex if and only if $\vertex{5}$ vanishes when one of its two-string arguments lies in that kernel; this is established in the finite-dimensional setting at the quintic order, its necessity needs a totally symmetric $\vertex{5}$ and an off-shell matching criterion, and the regularity of the kernel in infinite dimensions is open.\record{034-ker-v4-v5-compatibility (R38)}

\noindent\status{Proven} The nested kernels of the dressed tower depend on $\gs \Psi$ only and reduce at $\gs = 0$ to the kernels of the bare operators, and a level at which the kernel condition fails is visible in the value of the action; in the finite-dimensional setting, at fixed truncation.\record{061-ph2a-nested-kernels (R70)}

\noindent\status{Open} Whether the string vertices satisfy the kernel conditions for $k \ge 3$, in a suitable choice of the reduced vertices.\footnote{Ledger hypothesis PH2 (b).}

\subsubsection{Kernel conditions}
\label{pc:sec:all:kernel}

The auxiliary fields live on quotients, $\Sigma_k$ on $\Hn{k} / \ker \wprod{k}{k}$, and every term of the action must be a function on them.
At the quintic and septic orders this required the odd vertex to vanish on the kernel, and at the sextic order the source of the next field.
Besides its kinetic term, a kernel direction $N_k \in \ker \wprod{k}{k}$ enters three terms of the action: its own source, which pairs to zero with $N_k$ by the symmetry of $\wprod{k}{k}$, the odd coupling, and the source of $\Sigma_{k + 1}$.
The action descends to the quotients if the last two vanish on every kernel direction:
\begin{equation}
  \label{pc:eq:all:kernel-condition}
  \begin{gathered}
    \forall N_k \in \ker \wprod{k}{k}, \;
    \forall A_k \in \Hn{k}, \;
    \forall B \in \Hs:
    \\
    \dvertex{2 k + 1}(N_k, A_k, B)
      = 0,
    \qquad
    \wprod{k + 1}{k + 1}(B \otimes N_k)
      = 0,
  \end{gathered}
\end{equation}
for every $k \ge 2$.
The first clause is the condition of the odd orders and the second that of the even ones: at $k = 2$ the second clause is the vanishing of the sextic vertex on $\ker \sprod{2}{2}$,~\eqref{p:eq:hs6:kernel-condition}, since, for the form built from the diagrams, the exchange part of $\dvertex{6}$ vanishes there by the first clause.
The second clause says that a null direction of one level, tensored with any string field, is a null direction of the next: null directions propagate up the tower.
Every null vector of the kinetic operator of $\Sigma_k$ is a joint null vector of every vertex in which $\Sigma_k$ appears.
A kernel is meant in the square-integrable sense at fixed total momentum, as for $\sprod{2}{2}$ in \cref{p:sec:hs:quartic}.
Through the second clause, a null vector of $\sprod{2}{2}$ at a single total momentum would give the higher operators null vectors at every total momentum, but these are plane waves, sharp in one relative momentum, and do not count as kernel directions.
At $k = 2$ the condition bears on the quintic and sextic vertices themselves; for $k \ge 3$ it bears in part on the choice of the reduced vertices, which are defined up to terms vanishing at leading order.
The choice has a limit: at $k = 2$ a reduced sextic vertex satisfying the second clause exists only if its leading-order value vanishes on the symmetric tensors of $\Hn{3} \otimes \ker \sprod{2}{2} \otimes \Hs$.
The condition holds trivially when the bare operators have no kernel.
We do not try to settle here whether the string vertices satisfy it; in principle, this can be checked level by level in level truncation.

The conditions are what makes the quotient prescription consistent at every order.
The kernels in~\eqref{pc:eq:all:kernel-condition} are those of the bare operators, while the operators inverted in the elimination are the dressed ones, and the quotient is meaningful only if $\ker \wprod{k}{k}$ stays in the radical of the dressed operator for every $\Psi$, so that every dressed operator descends to the same quotient.
Read order by order in~\eqref{p:eq:hs8:scaling}, this is the condition: the first clause makes the term of order $\gs$ vanish on the kernel and the second the exchange term.
Generically the second clause is also necessary: a kernel direction it misses is lifted at order $\gs^2$.
When a clause fails along a kernel direction, that direction has no kinetic term at order $\gs^{2 k - 2}$ and acquires one from the odd coupling at the next order, or from the exchange one order further.
The failure does not show in the dressed operator of the level below at $\gs = 0$, so that the invertibility of the quartic form on its quotient does not certify the higher levels.
The construction as presented then breaks down and must be refined.\footnote{
  A way out may exist when a condition fails.
  The constraint $\Sigma_k = \Psi \otimes \Sigma_{k - 1}$ can be imposed by a multiplier instead of a kinetic term, either on the whole field, in which case nothing is inverted and no kernel condition arises, but the exchange structure is lost, or on the kernel directions only.
  This hybrid reproduces, at the sextic order, the original action through order $\gs^4$ for every complement of $\ker \sprod{2}{2}$, with a reduced vertex that depends on the complement (shown for a scalar model in \cref{p:app:scalar-hs}).
  The gauge transformations of the hybrid are left for future work.
}


\subsection{Signs of the composite terms and the order labels}
\label{ch3:sec:signs-labels}

Three later derivations complete what the blocks above say about the reduced vertices and the kinetic operators: the sign with which the exchange diagrams enter a reduced vertex, and which operators need the order label.
They use the assumptions of this chapter (tree level, formal power series in $\gs$, the kernel conditions) and the coefficients of the displays of the orders treated one at a time, which they take as given.

\paragraph{Sign of the composite terms.}

\noindent\status{Proven} In the basis of the diagrams of the cubic theory, the tree diagrams with $\Sigma_k$ propagators and the reduced vertices at their nodes, every composite term of a reduced vertex enters with a minus sign and a positive coefficient, at every order $n \ge 6$.\record{118-composite-vertex-sign (R114)}
By the rule~\eqref{p:eq:all:diagonal}, the leading-order value of the reduced vertex is the elementary vertex minus everything the cubic action of the previous order produces at order $\gs^{n - 2}$:
\begin{equation}
  \label{ch3:eq:composite-sign}
  \dvertex{n}(\Psi^n)
    = \vertex{n}(\Psi^n)
      - \sum_T n! \, c_T \, T,
  \qquad
  c_T
    > 0,
\end{equation}
where the sum runs over the tree diagrams $T$ of $\caction[\le n - 1]$ at order $\gs^{n - 2}$.
The coefficients are positive because the kinetic term of every auxiliary field carries a minus sign, which makes every propagator of the eliminated action positive, while the odd couplings carry a plus sign; the sources enter linearly and twice per chain of $\Sigma_k$ lines, at its two ends, so that their sign drops out.
The reality signs of the auxiliary fields and the exchange sign of the pairing are conditions on the fields, not coefficients of the action, and do not enter.
Distinct diagrams are taken to evaluate to independent forms, as the septic block assumes for the string vertices; this is what makes \enquote{the coefficient of a diagram} well defined.

Re-expanded in the elementary string vertices, the sign of a composite term depends on the diagram, because a diagram nested in a reduced node is counted once more with the opposite sign.
The first positive elementary term is the chain of three quintic vertices in $\dvertex{7}$, written $\mathcal S_1$ in the septic block, whose coefficient in~\eqref{p:eq:hs7:diagonal} is
\begin{equation}
  \label{ch3:eq:chain-sign}
  + \frac{42}{25}
    = - \frac{42}{25}
      + \frac{14}{5} \cdot \frac{6}{5}:
\end{equation}
the chain is subtracted once directly and added back through the diagram $\dvertex{6} \cdot \vertex{5}$, whose reduced node contains it with the coefficient $- \tfrac{6}{5}$; the sextic--quintic diagram $\mathcal S_2$ keeps the coefficient $- \tfrac{14}{5}$.
The rational identities and the positivity of every monomial of the stationary value in a two-level model, to third order in the coupling, are machine-checked, and the check fails when the sign of the septic vertex is flipped.

In words: the reduced vertex always subtracts the exchanges of the auxiliary fields, as long as these exchanges are written with the vertices of the cubic theory at their nodes.
This is what supports the reading of each exchange diagram as a region removed from the vertex region of $\mathcal V_{0, n}$, in the covering question of \cref{ch:worldsheet-circle}.
The statement concerns the sign of the coefficients only; it does not say whether the exchanges are integrals over regions of moduli space.

\paragraph{Order labels of the kinetic operators.}

\noindent\status{no ledger entry} The quadratic form of $\Sigma_k$ that does not depend on $\Psi$ is $\tfrac{\gs^{2 k - 2}}{(2 k)!} \, \wprod{k}{k}$ at every order in $\gs$, for every truncation $K \ge k$ of the tower and after the elimination of any subset of the other auxiliary fields, and $\wprod{k}{k}$ does not depend on $K$.\record{115-bare-operator-notation}
At $\Psi = 0$ the auxiliary sector is block-diagonal in the tower, so that eliminating any auxiliary field adds only $\Psi$-dependent terms to the quadratic forms of the others, and the reduced even vertex $\dvertex{2 k}$ is fixed at the order $\gs^{2 k - 2}$ at which $\Sigma_k$ enters.
The bare operator therefore needs no order label, which is why the later blocks write $\wprod{k}{k}$ where the earlier ones write $\kin{k}{}{2 k - 2}$.

\noindent\status{no ledger entry} The dressed operators do need the label: they receive corrections that depend on $\Psi$ alone and are not linear in it, from the elimination of the higher auxiliary fields, from relative order $\gs^2$ on.\record{117-dressed-operator-order-label}
For $\Sigma_2$, the elimination of $\Sigma_3$ at the sextic order adds a term quadratic in $\Psi$, exactly in $\gs$,
\begin{equation}
  \label{ch3:eq:label-k2}
  \pair[2]{A}{\kin{2}{\Psi}{4}(B)}
    - \pair[2]{A}{\kin{2}{\Psi}{3}(B)}
    = - \frac{\gs^2}{30} \, \dvertex{6}(\Psi \otimes A, \Psi \otimes B),
\end{equation}
the septic truncation adds terms of degree three and four in $\Psi$ at the orders $\gs^3$ and $\gs^4$, and the octic truncation a further term of degree four, $- \tfrac{\gs^4}{1680} \, \dvertex{8}(\Psi^2 \otimes A, \Psi^2 \otimes B) + O(\gs^5)$.
In general, with the superscript $n$ the power of $\gs$ at which the action is truncated, $n \ge 2 k - 2$, the truncated operator agrees with the operator of the full tower up to $O(\gs^{n - 2 k + 3})$, with a leading difference that does not vanish in general.
The rule the blocks follow is thus: $\kin{k}{\Psi}{n}$ is the operator of the action truncated at $\gs^n$, and the operator of the full tower carries no superscript.
Both statements are checked in a finite-dimensional model by companion scripts, whose break legs fail as they should.

\subsection{Two complements from the derivation records}
\label{ch3:sec:complements}

\noindent\status{no ledger entry} Both complements are pure algebra on top of the equations above, and both are checked in a finite-dimensional model by companion scripts.\record{100-all-orders-action-eom}\textsuperscript{,}\record{085-paper-gaps-3-4}

\subsubsection{First orders of the solution through \texorpdfstring{$\Sigma_3$}{Sigma3}}
\label{ch3:sec:first-orders}

The order-by-order recursion~\eqref{p:eq:all:eom-orders} determines the solution of the auxiliary equations, but it is useful to see its first terms written out, since they show how a correction travels along the tower.
We take the septic truncation, with the fields $\Psi$, $\Sigma_2$, $\Sigma_3$ and the action $\caction[\le 7]$ of the septic order, which is the series~\eqref{p:eq:all:action} truncated at $\gs^5$.
The inverses below are those of the bare operators on their quotients, so that the solutions are defined modulo the bare kernels.

\paragraph{The equations.}

At $K = 3$ the general equations~\eqref{p:eq:all:eom-psi} and~\eqref{p:eq:all:eom-sigma} become:
\begin{subequations}
  \label{ch3:eq:eom-K3}
  \begin{align}
    \label{ch3:eq:eom-K3-psi}
    \mathcal E_\Psi
      &
      = Q_B \Psi
        + \frac{\gs}{2} \, \ell_2(\Psi^2)
        + \frac{\gs^2}{3!} \, \ell_3(\Psi, \Sigma_2)
        + \frac{\gs^3}{5!} \, \ell_4(\Sigma_2^2)
      \notag
      \\ &
      \quad
        + \frac{2 \gs^4}{6!} \, \dvertex{6}(\, \cdot \, \otimes \Sigma_2, \Sigma_3)
        + \frac{\gs^5}{7!} \, \dvertex{7}(\Sigma_3, \Sigma_3, \, \cdot \,),
    \\
    \label{ch3:eq:eom-K3-sigma2}
    \mathcal E_{\Sigma_2}
      &
      = \sprod{2}{2}(\Psi^2)
        - \sprod{2}{2}(\Sigma_2)
        + \frac{\gs}{5} \, \sprod{3}{2}(\Psi, \Sigma_2)
        + \frac{\gs^2}{30} \, \dvertex{6}(\Psi \otimes \, \cdot \,, \Sigma_3),
    \\
    \label{ch3:eq:eom-K3-sigma3}
    \mathcal E_{\Sigma_3}
      &
      = \wprod{3}{3}(\Psi \otimes \Sigma_2)
        - \wprod{3}{3}(\Sigma_3)
        + \frac{\gs}{7} \, \dvertex{7}(\, \cdot \,, \Sigma_3, \Psi)
      = \wprod{3}{3}(\Psi \otimes \Sigma_2)
        - \kin{3}{\Psi}{5}(\Sigma_3).
  \end{align}
\end{subequations}
The first is the $\Psi$-equation of the sextic order with the septic term added, the second is the $\Sigma_2$-equation of the sextic order, and the third is the $\Sigma_3$-equation~\eqref{p:eq:hs7:saddle-3} of the septic order.
In the second, the coefficient $\gs^2 / 30 = \gs^2 / (5 \cdot 6)$ is the factor $\gs^2 / ((2 k + 1) (2 k + 2))$ of~\eqref{p:eq:all:eom-sigma} at $k = 2$.
At $k = 3 = K$ there is no heavier field, and the exchange term is absent from the third equation.

\paragraph{The solutions.}

We expand $\Sigma_k = \sum_m \gs^m \Sigma_k^{[m]}$ with $\Sigma_k^{[0]} = \Psi^k$, as in~\eqref{p:eq:all:eom-orders}.
At order $\gs$, the recursion for $k = 2$ has no correction from the lighter field, since $\Sigma_1 = \Psi$ is not expanded, and no exchange term, which starts at order $\gs^2$; only the quintic term acts:
\begin{equation}
  \label{ch3:eq:sigma2-1}
  \Sigma_2^{[1]}
    = \frac{1}{5} \, \sprod{2}{2}^{-1} \circ \sprod{3}{2}(\Psi, \Psi^2).
\end{equation}
For $k = 3$ at the same order, the source carries the correction of $\Sigma_2$ and the septic vertex acts on the leading term $\Psi^3$:
\begin{equation}
  \label{ch3:eq:sigma3-1}
  \Sigma_3^{[1]}
    = \Psi \otimes \Sigma_2^{[1]}
      + \frac{1}{7} \, \wprod{3}{3}^{-1} \circ \dvertex{7}(\, \cdot \,, \Psi^3, \Psi).
\end{equation}
Together:
\begin{equation}
  \label{ch3:eq:first-orders}
  \begin{aligned}
    \Sigma_2^*
      &
      = \Psi^2
        + \frac{\gs}{5} \, \sprod{2}{2}^{-1} \circ \sprod{3}{2}(\Psi, \Psi^2)
        + O(\gs^2),
    \\
    \Sigma_3^*
      &
      = \Psi^3
        + \gs \Big[
          \Psi \otimes \frac{1}{5} \, \sprod{2}{2}^{-1} \circ \sprod{3}{2}(\Psi, \Psi^2)
          + \frac{1}{7} \, \wprod{3}{3}^{-1} \circ \dvertex{7}(\, \cdot \,, \Psi^3, \Psi)
        \Big]
        + O(\gs^2).
  \end{aligned}
\end{equation}
The first line is the expansion of the quintic solution of \cref{ch:order-by-order}.
At order $\gs^2$, the recursion for $k = 2$ receives two terms, the quintic term acting on $\Sigma_2^{[1]}$ and the exchange of $\Sigma_3$ evaluated at its leading term:
\begin{equation}
  \label{ch3:eq:sigma2-2}
  \Sigma_2^{[2]}
    = \frac{1}{25} \, \Big( \sprod{2}{2}^{-1} \circ \sprod{3}{2}(\Psi, \, \cdot \,) \Big)^2 \Psi^2
      + \frac{1}{30} \, \sprod{2}{2}^{-1} \circ \dvertex{6}(\Psi \otimes \, \cdot \,, \Psi^3).
\end{equation}
The first term is the iterate of the quintic dressing, the second the $\Sigma_3$ exchange.

\paragraph{What the first orders show.}

The order-$\gs$ correction reaches $\Sigma_3$ by two routes: through its source, which carries the correction of the lighter field $\Sigma_2$ at the same order, and through its own odd vertex $\dvertex{7}$ acting on the leading term.
The heavier field acts on the lighter one only through the exchange term, two orders later: the $\Sigma_3$ exchange is present in the $\Sigma_2$-equation from the sextic order on, and enters $\Sigma_2^*$ at order $\gs^2$, while $\dvertex{7}$ reaches $\Sigma_2^*$ only through $\Sigma_3^*$, at order $\gs^3$.
This is the triangular structure of the recursion~\eqref{p:eq:all:eom-orders} in its smallest non-trivial case: corrections move up the tower at the same order, act on their own level one order later through the odd vertex, and move down the tower two orders later.

\subsubsection{The converse of the kernel-shift statement}
\label{ch3:sec:converse}

The recovered-action blocks show that, under the kernel condition~\eqref{pc:eq:all:kernel-condition}, a shift of the auxiliary fields along the bare kernels changes neither the equations nor the value of the action.
The converse also holds: asking only that the equation of motion of $\Psi$ be unchanged by every kernel shift, at arbitrary values of the auxiliary fields, already forces both clauses of the condition, in its field-space form (at the top level of a truncated tower, only the first clause).
The argument is pure algebra and does not use the matching theorem.
We work at a fixed level truncation with the momenta fixed or discretized, so that every space is finite-dimensional, with the tower truncated at $K$, the bare operators symmetric for the $k$-string pairing and the pairings nondegenerate.

\paragraph{The shift of the action.}

Write $a_k := \gs^{2 k - 2} / (2 k)!$ and $b_k := \gs^{2 k - 1} / (2 k + 1)!$, so that the action~\eqref{p:eq:all:action} truncated at $K$ is
\begin{equation}
  \label{ch3:eq:action-abc}
  \begin{aligned}
    \caction[\le 2 K + 1]
      &
      = S_{\le 3}
        + \sum_{k = 2}^{K} \Big[
          - a_k \, \pair[k]{\Sigma_k}{\wprod{k}{k}(\Sigma_k)}
          + 2 a_k \, \pair[k]{\Sigma_k}{\wprod{k}{k}(\Psi \otimes \Sigma_{k - 1})}
      \\ &
      \qquad\qquad\qquad
          + b_k \, \dvertex{2 k + 1}(\Sigma_k, \Sigma_k, \Psi)
        \Big],
  \end{aligned}
\end{equation}
with $\Sigma_1 := \Psi$.
Shift $\Sigma_k \to \Sigma_k + N_k$ for $k = 2, \ldots, K$, with arbitrary $N_k \in \ker \wprod{k}{k}$, and leave $\Psi$ unshifted ($N_1 := 0$).
By the symmetry of $\wprod{k}{k}$ and $\wprod{k}{k}(N_k) = 0$, the kinetic term of each level does not change, and the source of level $k$ changes only through the shift of its lighter field:
\begin{equation}
  \label{ch3:eq:shift-terms}
  \begin{aligned}
    \Delta \pair[k]{\Sigma_k}{\wprod{k}{k}(\Sigma_k)}
      &
      = 2 \, \pair[k]{N_k}{\wprod{k}{k}(\Sigma_k)}
        + \pair[k]{N_k}{\wprod{k}{k}(N_k)}
      = 0,
    \\
    \Delta \pair[k]{\Sigma_k}{\wprod{k}{k}(\Psi \otimes \Sigma_{k - 1})}
      &
      = \pair[k]{\Sigma_k}{\wprod{k}{k}(\Psi \otimes N_{k - 1})},
    \\
    \Delta \, \dvertex{2 k + 1}(\Sigma_k, \Sigma_k, \Psi)
      &
      = 2 \, \dvertex{2 k + 1}(N_k, \Sigma_k, \Psi)
        + \dvertex{2 k + 1}(N_k, N_k, \Psi).
  \end{aligned}
\end{equation}
In the second line, the part with $N_k$ in the slot of $\Sigma_k$ vanishes by symmetry, whatever the other argument.
The total shift of the action, $D(\Psi; \Sigma, N) := \caction(\Psi, \Sigma + N) - \caction(\Psi, \Sigma)$, is then:
\begin{equation}
  \label{ch3:eq:D}
  \begin{aligned}
    D(\Psi; \Sigma, N)
      &
      = \sum_{k = 2}^{K - 1} 2 a_{k + 1} \, \pair[k + 1]{\Sigma_{k + 1}}{\wprod{k + 1}{k + 1}(\Psi \otimes N_k)}
      \\ &
      \quad
        + \sum_{k = 2}^{K} b_k \Big[
          2 \, \dvertex{2 k + 1}(N_k, \Sigma_k, \Psi)
          + \dvertex{2 k + 1}(N_k, N_k, \Psi)
        \Big].
  \end{aligned}
\end{equation}
Every term of $D$ is linear in $\Psi$.
The only term of the action quadratic in $\Psi$, the source of $\Sigma_2$, has dropped out: its shift is the pairing of $N_2$ with $\sprod{2}{2}(\Psi^2)$, which vanishes by the symmetry of $\sprod{2}{2}$, and $\Psi$ is not shifted.
For a single level, \eqref{ch3:eq:D} is the shift~\eqref{pc:eq:all:kernel-shift}.

\paragraph{The equivalence.}

Since $N$ does not depend on $\Psi$, the change of the $\Psi$-equation under the shift is the derivative of $D$:
\begin{equation}
  \label{ch3:eq:delta-E}
  \mathcal E_\Psi(\Psi, \Sigma + N) - \mathcal E_\Psi(\Psi, \Sigma)
    = \frac{\delta D}{\delta \Psi},
\end{equation}
and the change of the equation of $\Sigma_k$ is, up to its normalization in~\eqref{p:eq:all:eom}, the derivative $\delta D / \delta \Sigma_k$.
If the kernel condition holds, every term of~\eqref{ch3:eq:D} vanishes, the first clause with the arguments $(\Sigma_k, \Psi)$ and $(N_k, \Psi)$, the second with $B = \Psi$; then $D \equiv 0$ and all the equations are the same at $\Sigma$ and at $\Sigma + N$.
This is the direction used in the paper.

Conversely, suppose that $\mathcal E_\Psi(\Psi, \Sigma + N) = \mathcal E_\Psi(\Psi, \Sigma)$ for all $\Psi$, $\Sigma$ and $N$.
Then $\delta D / \delta \Psi = 0$, and since $D$ is homogeneous of degree one in $\Psi$, Euler's relation gives:
\begin{equation}
  \label{ch3:eq:euler}
  D(\Psi; \Sigma, N)
    = \pair{\Psi}{\frac{\delta D}{\delta \Psi}}
    = 0.
\end{equation}
Keep a single non-zero $N_k$ and rescale it, $N_k \to \lambda N_k$.
The terms of order $\lambda^2$ and $\lambda$ of $D$ vanish separately:
\begin{equation}
  \label{ch3:eq:lambda}
  \begin{gathered}
    \dvertex{2 k + 1}(N_k, N_k, \Psi)
      = 0,
    \\
    2 a_{k + 1} \, \pair[k + 1]{\Sigma_{k + 1}}{\wprod{k + 1}{k + 1}(\Psi \otimes N_k)}
      + 2 b_k \, \dvertex{2 k + 1}(N_k, \Sigma_k, \Psi)
      = 0.
  \end{gathered}
\end{equation}
Since $\Sigma_k$, $\Sigma_{k + 1}$ and $N_k$ are independent variables, we can set $\Sigma_{k + 1} = 0$ in the second relation, which gives $\dvertex{2 k + 1}(N_k, A, \Psi) = 0$ for every $A$, the first clause, for $k \le K$.
Setting $\Sigma_k = 0$ instead and using the nondegeneracy of the $(k + 1)$-string pairing gives $\wprod{k + 1}{k + 1}(\Psi \otimes N_k) = 0$, the second clause, for $k \le K - 1$.
Both clauses are linear in the string field, which ranges over the field space, so they hold for every $B$ in its span.
Hence:
\begin{equation}
  \label{ch3:eq:converse}
  \begin{gathered}
    \Big(
      \mathcal E_\Psi(\Psi, \Sigma + N)
        = \mathcal E_\Psi(\Psi, \Sigma)
      \ \ \forall \, \Psi, \Sigma, N
    \Big)
    \iff
    D \equiv 0
    \\
    \iff
    \text{the kernel condition, in its field-space form}.
  \end{gathered}
\end{equation}

\paragraph{Scope.}

The converse is an off-shell statement.
On the solution, $\Sigma_k^*$ and $\Sigma_{k + 1}^*$ are fixed functions of $\Psi$, and the order-$\lambda$ relation of~\eqref{ch3:eq:lambda} gives only one combination of the two clauses; moreover, without the condition, $\Sigma^* + N$ is in general not a solution, so that the set of representatives of the solution is not defined.
At the top level $k = K$ there is no source of a heavier field, and only the first clause is obtained; in the even truncation the top odd coupling is absent, and the first clause at $k = K$ is not needed.
The clauses obtained hold with $A$, $B$ and $N_k$ in the spaces of the fields, of ghost numbers $2 k$, $2$ and $2 k$.
The condition~\eqref{pc:eq:all:kernel-condition}, written for all $A \in \Hn{k}$ and $B \in \Hs$, is a stronger sufficient condition: for $B$ of ghost number other than $2$, the term $\wprod{k + 1}{k + 1}(B \otimes N_k)$ is not removed by ghost-number counting.

\paragraph{What the converse means.}

Every term through which a kernel direction could reach the action is linear in $\Psi$, so the equation of motion of $\Psi$ alone already sees all of them.
The kernel condition is therefore not only sufficient for the quotient prescription but forced, as soon as one asks that the equation of the string field not depend on how the auxiliary fields are represented on the quotients, off shell.
The statement does not say whether the string vertices satisfy the condition, which remains open for $k \ge 3$.

\subsection{Three bookkeepings and the fixed-point theorem}
\label{ch3:sec:schemes}

The last two blocks come from the earlier pedagogical notes on the construction and its algebraic formulation.
They compare three ways of organizing the subtractions that the all-orders problem requires, and then prove the all-orders theorem in the bookkeeping that the paper adopts.
They use the notation of those notes, which differs from the rest of this document as follows.

\begin{quote}
  \small
  \textbf{Translation box.}
  \begin{itemize}
    \setlength{\itemsep}{0.3ex}

    \item $V_n$ is the elementary vertex $\vertex{n}$; $\tfrac12 V_2 + \tfrac{\gs}{3!} V_3$ is $S_{\le 3}$, the kinetic term and the cubic vertex of $\Psi$.

    \item $V_4(\Sig, \Sig)$ as the kinetic term of $\Sig_2$ is $\vertex{4}(\Sigma_2, \Sigma_2) = \pair[2]{\Sigma_2}{\sprod{2}{2}(\Sigma_2)}$; $V_4^{-1}$ is $\sprod{2}{2}^{-1}$ on the quotient.

    \item $\Sym^k$ is the $k$-string space $\Hn{k}$, and $P_k := \Psi^{\otimes k}$ (the Veronese element) is $\Psi^k$.

    \item Scheme B (dressed, uniform) is the all-order action~\eqref{p:eq:all:action} of this chapter: its running effective vertices $V_{n, \rm eff}$ and $\tilde V_n$ are the reduced vertices $\dvertex{n}$; $\tilde M_k$ is the bare kinetic operator $\wprod{k}{k} = \kin{k}{}{2 k - 2}$; the source $\Mean{\tilde u_k(\Sig_{k - 1}), \Sig_k}$ is $\pair[k]{\Sigma_k}{\wprod{k}{k}(\Psi \otimes \Sigma_{k - 1})}$; $\hat M_k$ is the operator $\kin{k}{\Psi}{2 k - 1}$ dressed by the odd vertex~\eqref{p:eq:all:dressed-odd}; $L_k(\Psi)$ is the source map $\wprod{k}{k}(\Psi \otimes \, \cdot \,)$; $\hat H_k$ is the nested operator~\eqref{p:eq:all:nested}; the tensors $T_n$ of the theorem are the reduced vertices, whose diagonals the recursion~\eqref{ped:eq:diagrec} fixes as~\eqref{p:eq:all:diagonal} does.

    \item Scheme A (minimal) is the presentation with the fewest auxiliary fields, $\Sigma_2$ alone through the sextic order; its $W_6$ is the reduced sextic vertex of the form whose sextic term is written with $\Sigma_2$ alone (\cref{ch:order-by-order}).

    \item Scheme C (elementary, indefinite) is the partner form, and the wrong-sign partner $\chi$ is $\Upsilon_2$ (\cref{ch:partner-field}).

    \item Route B is the order-by-order construction of \cref{ch:order-by-order}, which cubicizes the action one order at a time.

    \item Reality: the notes impose $\Sig^\ddagger = \Sig$, which is $\Sigma_k^\star = (-1)^k \Sigma_k$ in the star convention of this document (\cref{ch:conventions}), since $\ddagger = (-1)^k \star$ on a $k$-string field; a reader who reads $\ddagger$ as $\star$ gets the sign wrong for odd $k$, in particular for $\Sigma_3$.
      The condition does not appear in these two blocks.
  \end{itemize}
\end{quote}

The first block describes the three schemes.
The results it quotes for the individual orders are stated with their status where they are derived: the quartic to septic orders in \cref{ch:order-by-order}, the wrong-sign partner in \cref{ch:partner-field}.
The partner bookkeeping beyond the order $\gs^4$, which the block describes as open, is taken up in \cref{ch:partner-field}.

\subsubsection{Scheme A (minimal): composite vertices \texorpdfstring{$W_n$}{Wn} on the
smallest tower}
\label{ped:sec:allorders:min}

Use the fewest fields ($\Sig_k$ first at $n=3(k{-}1)+1$, i.e.\ $k\le\lceil n/3\rceil$,
by the all-orders counting) and put every subtraction \emph{inside the vertex kernel}: at order $n$
add the single coupling
\begin{equation}
	\frac{\gs^{n-2}}{n!}\,W_n(\Sig_a,\Sig_b,\Sig_c),
	\qquad n=a+b+c,\quad\text{balanced partition } a,b,c\le\lceil n/3\rceil,
	\label{ped:eq:schemeA}
\end{equation}
with the recursion established through $n=6$:
\begin{equation}
	W_4=V_4,\qquad W_5=V_5,\qquad
	W_6=V_6-\tfrac65\,\big(W_5\circ V_4^{-1}\circ W_5\big),\qquad\dots
	\label{ped:eq:Wnrecursion}
\end{equation}
--- in general $W_n=V_n-(\text{all colour-2-reducible composites of lower }W_m)$, the
coloured-operad statement that the coloured vertex equals the amplitude minus the exchange
regions, with the $\Sig$-propagators moduli-free kernels. A reading
of why this scheme is natural: the subtraction removes from the
elementary vertex region exactly the part that can be \emph{covered} by lower-order
auxiliary-field exchanges --- mirroring how the original closed-SFT action is itself
constructed, by identifying the gaps in the moduli-space covering left by propagator
regions and covering them with vertex regions; here the \enquote{propagators} are the
moduli-free colour-2 necks and the covering bookkeeping happens at the level of kernels
rather than regions. Proven:
$n\le6$ (\cref{p:sec:hs:quintic,p:sec:hs:sextic}). At $n=7$ the literal one-coupling rule
\eqref{ped:eq:schemeA} is now \emph{falsified} (\cref{p:sec:hs:septic}):
$\Sig_3$ becomes unavoidable ($7=3+2+2$ is the only partition reaching seven strands),
but a $(3,2,2)$ coupling alone cannot supply the $V_7$ contact --- with no $O(\gs^4)$
source the $\Sig_3$ saddle is $O(\gs)$ and the coupling's on-shell footprint is
$O(\gs^6)$, for \emph{any} kernel and coefficient. The $\Sig_3$
Hubbard--Stratonovich pair at $O(\gs^4)$ (scheme B's even pair) is \emph{forced}, so at
$n=7$ scheme A must restructure its $n=6$ sector and merge into the scheme-B family
(\cref{p:sec:hs:septic}).

\subsubsection{Scheme B (dressed, uniform): pair each field with the running effective
vertex}
\label{ped:sec:allorders:dressed}

The preferred bookkeeping (its first instance, at the sextic order, is proven,
\cref{p:sec:hs:sextic}). Define the \emph{running effective vertex}
\begin{equation}
	V_{n,\rm eff}
	:= V_n+\big(\text{all feedback composites generated by the orders}<n\big)
	\label{ped:eq:runningV}
\end{equation}
--- operationally, $V_{n,\rm eff}$ is the $n$-th derivative, on the Veronese variety, of
the action obtained by eliminating the auxiliaries of the orders $<n$; at $n=6$ this is
$\tilde V_6$ with $\tilde V_6(P_3,P_3)=W_6(\Psi^6)$, built by the $\Phi$-map
\eqref{ped:eq:tildeMdef}. The scheme then adds, for each $k\ge2$, one uniform triple of
terms --- an HS pair (kinetic $+$ source, ratio $1{:}2$) at the even order $n=2k$ and
one odd-order coupling at $n=2k+1$:
\begin{equation}
	\boxed{\;
	S^{\rm dr}
	= \tfrac12V_2(\Psi^2)+\frac{\gs}{3!}V_3(\Psi^3)
	+\sum_{k\ge2}\bigg[
	\frac{\gs^{2k-2}}{(2k)!}\Big(-\tilde M_k(\Sig_k,\Sig_k)
	+2\,\Mean{\tilde u_k(\Sig_{k-1}),\Sig_k}\Big)
	+\frac{\gs^{2k-1}}{(2k+1)!}\,\tilde V_{2k+1}(\Sig_k,\Sig_k,\Psi)
	\bigg],
	\;}
	\label{ped:eq:Sdrallorders}
\end{equation}
with $\tilde M_k:=V_{2k,\rm eff}\big|_{\Sym^k\times\Sym^k}$ and
$\Mean{\tilde u_k(\Sig_{k-1}),\Sig_k}:=\tilde V_{2k,\rm eff}(\Sig_k,\Sig_{k-1},\Psi)$
(partitions $2k=k+k$ and $2k=k+(k{-}1)+1$; every term has at most three fields). The
$k=2$ triple is exactly \eqref{p:eq:hs4:action-tilde} plus the quintic coupling ($\tilde M_2=V_4$,
$\tilde V_5=V_5$: no dressing yet); the $k=3$ even pair is \eqref{ped:eq:Sdr}
(proven); the $k=3$ odd coupling $\tilde V_7(\Sig_3,\Sig_3,\Psi)$ at
$O(\gs^5)$ is now \emph{proven too} (\cref{p:sec:hs:septic}), with the
explicit forced dressing \eqref{p:eq:hs7:dressed-vertex} --- the running rule \eqref{ped:eq:runningV}
verbatim at $n=7$; and the whole tower is now \emph{proven at every order}
(\cref{ped:sec:allorders:theorem}: the action is jointly quadratic in the
tower, the effective action is one Gaussian with an operator-continued-fraction
kernel, and closure is a solvable forced-diagonal recursion). What is uniform and established across $k$: the $1{:}2$ ratio, the
Veronese saddle $\Sig_k^*=P_k$ exactly on-shell (the collapse identities
\eqref{ped:eq:tildecollapse} pattern), homogeneous saddle-tracking gauge laws with no new
parameters, and $\IBL$ products dressed by lower
composites satisfying no new relations (\cref{ch:ibl}). The vertices grow composite as the dressing accumulates,
but the \emph{rule} \eqref{ped:eq:runningV} is one line.

\subsubsection{Scheme C (elementary, indefinite): geometric vertices plus wrong-sign
partners}
\label{ped:sec:allorders:chi}

The $\chi$-mechanism of \cref{ped:sec:milestones:chi} suggests the third bookkeeping: keep
\emph{every} vertex elementary/geometric and delegate \emph{all} subtractions to
wrong-sign partner fields. Established at $O(\gs^4)$: \eqref{ped:eq:S4chi}.
The underlying principle is general --- every composite that Route B
ever needs is the induced action of a linearly-sourced Gaussian (the Liouville
mechanism), hence can be re-localised by a partner field with the opposite-sign kernel;
what is open is the bookkeeping of the partners at $O(\gs^5)$ and beyond (the $O(\gs^5)$
feedback is now explicit, \eqref{p:eq:hs7:induced}: two strand-local chains $S_1,S_2$ in the
dressed presentation, plus the strand-split $S_3$ in the minimal one --- so either a
single $\chi$ with new couplings or a wrong-sign tower $\chi_k$ mirroring $\Sig_k$).
If it closes, this scheme is the
most geometric cubic form: all kernels are the geometric $V_n$ themselves, at
the price of an indefinite (split-signature) auxiliary sector --- which
\eqref{ped:eq:lightcone} suggests is not a bug but the signature of an enlarged gauge-fixed
theory.



\medskip
The second block proves the all-orders theorem in scheme B: the Jacobi form of the action in the tower, the continued-fraction kernel, the telescoping contact lemma, the theorem, the gauge laws at all orders, and the instance at order $\gs^6$ checked in a finite-dimensional model.
Its continued fraction~\eqref{ped:eq:contfrac} is the nested operator~\eqref{p:eq:all:nested}, its telescoping contact lemma corresponds to the collapse lemma of the first block together with its statement that the data of level $k$ first enter the value of the action at the order at which they are introduced, and its gauge law~\eqref{ped:eq:dYall} is the law of the subsection on gauge transformations at all orders.

\noindent\status{Conditional(PH1, PH2 (b), PH5, PH14 [through $A_6$]) + R70 -- Proven} The all-orders fixed-point theorem: within the jointly Gaussian class with a shared parent form for the kinetic term and the source of each level, a cubic action matching the original one exists at every order as a formal power series in $\gs$, with the diagonals of the reduced vertices forced and everything else presentation freedom.
The scope stated at the end of the block binds: couplings of $(3, 3, 2)$ type at $n \ge 8$, the convergence of the continued fraction and the analytic properties of the nested operators for the string vertices are not covered.\record{019-b7-allorders-fixed-point (R26)}

\subsubsection{The all-orders theorem: the fixed point solved}
\label{ped:sec:allorders:theorem}

The fixed-point proof announced above has been carried out (machine-checked end to end
at truncation $K=4$ by a companion script, listed in the script table of
\cref{ch:finite-models}: 8/8 checks on generic
non-$\Loo$ tensors $V_4\dots V_9$, two seeds). Three structural facts, then the
theorem.

\subparagraph{(i) The scheme-B action is a Jacobi quadratic form in the tower.} Stronger
than sequential Gaussianity: \eqref{ped:eq:Sdrallorders} is \emph{jointly} quadratic in
$(\Sig_2,\dots,\Sig_K)$. With $a_k:=\gs^{2k-2}/(2k)!$ its Hessian is
block-\emph{tridiagonal} --- diagonal blocks $a_k\hat M_k$ with
\begin{equation*}
	\hat M_k=\tilde M_k-\tfrac{\gs}{2k+1}\,\tilde V_{2k+1}(\cdot,\cdot,\Psi),\qquad
	L_k(\Psi)=\tilde V_{2k}(\cdot,\cdot,\Psi),
\end{equation*}
(the odd couplings are kinetic dressings, never sources, because their third slot is
$\Psi$); off-diagonal blocks the source maps $L_k(\Psi)$ (partition
$k+(k{-}1)+1$); and the only inhomogeneity is the $k=2$ source
$\tilde u_2=V_4(\cdot,\Psi^{\otimes2})$. Eliminating the tower is therefore \emph{one}
exact Gaussian, and the on-shell value is the $(2,2)$ resolvent block of a Jacobi
operator --- an operator continued fraction:
\begin{equation}
	\boxed{\;
	\begin{aligned}
	S(\text{saddles})
	&= \tfrac12V_2+\tfrac{\gs}{3!}V_3
	+ \frac{\gs^2}{4!}\,\tilde u_2^{\top}\hat H_2(\gs,\Psi)^{-1}\tilde u_2,\\
	\hat H_k &= \hat M_k - \gs^2\,\tfrac{(2k)!}{(2k+2)!}\,
	L_{k+1}^{\top}\hat H_{k+1}^{-1}L_{k+1}.
	\end{aligned}
	\;}
	\label{ped:eq:contfrac}
\end{equation}
The $k=2$ instance of the recursion is the dressed operator~\eqref{p:eq:hs7:dressed-2} of
the septic order (the factor $(2k)!/(2k+2)!$ gives its $\gs^2/30$ at $k=2$, then
$\gs^2/56$ at $k=3$). Closure
\eqref{p:eq:all:closure} becomes the statement that the \emph{full} classical effective
action is the value of a single Gaussian in $\Sig_2$ with the continued-fraction
kernel.

\subparagraph{(ii) Telescoping contact lemma.} Level-$k$ data first enter the
$\gs$-expansion of \eqref{ped:eq:contfrac} at order $\gs^{2k-2}$ (each rung of the
fraction costs $\gs^2$), and at the two leading orders $\gs^{2k-2},\gs^{2k-1}$ they
enter \emph{only through the Veronese diagonals} $\tilde V_{2k}(\Psi^{2k})$,
$\tilde V_{2k+1}(\Psi^{2k+1})$, linearly, with coefficients exactly $1/(2k)!$ and
$1/(2k+1)!$. Two ingredients: the collapse identity $\tilde u_j(P_{j-1})=\tilde M_j
P_j$ pins every flank of the insertion to the Veronese at leading order (it holds
automatically whenever kinetic and source descend from one parent form --- block
composites like $\Phi^\top V_4^{-1}\Phi$ and fully symmetric polarizations equally);
and the rung prefactors telescope,
$\tfrac1{4!}\cdot\tfrac{4!}{6!}\tfrac{6!}{8!}\cdots\tfrac{(2k-2)!}{(2k)!}
=\tfrac1{(2k)!}$. (Machine form: perturbing the level-4 parent $T_8$ by
$\epsilon\,a^{\otimes8}$ shifts the $\gs^6$ coefficient by exactly
$\epsilon\,(a\!\cdot\!\Psi)^8/8!$ and leaves the lower coefficients bit-identical.)

\subparagraph{(iii) Theorem (all-orders closure).} Let the $V_n$ be arbitrary symmetric
tensors with the nested kernels $\hat H_k(\gs,\Psi)$ invertible near $\gs=0$ (the
all-orders auxiliarity hypothesis; for the string vertices, the analytic properties of
the nested kernels are not addressed here). Define
$T_4=V_4$, $T_5=V_5$ and recursively
\begin{equation}
	T_n := -\,n!\,\mathrm{Pol}\big(A_{n-2}\big),\qquad
	A_{n-2}(\Psi) := [\gs^{n-2}]\big(S_{<n}(\text{saddles})-S_{\Loo}\big),
	\label{ped:eq:diagrec}
\end{equation}
where $S_{<n}$ is the tower built from $T_4,\dots,T_{n-1}$ and $\mathrm{Pol}$ is
polarization ($A_{n-2}$ is a homogeneous degree-$n$ polynomial in $\Psi$, by the
scaling $\Psi\to t\Psi$, $\Sig_k\to t^k\Sig_k$, $\gs\to\gs/t$ under which every term
of the action scales as $t^2$ --- so a symmetric tensor with that diagonal always
exists; the geometric $V_n$ reappears through $\mathrm{Pol}$, i.e.\
$T_n=V_n-n!\,\mathrm{Pol}(\text{feedback})$, the running rule \eqref{ped:eq:runningV}
verbatim). Then \eqref{ped:eq:Sdrallorders} with these kernels satisfies
\eqref{p:eq:all:closure} \emph{as a formal power series in $\gs$}, and at every finite
truncation $K$ the identity is exact-rational, matched through $O(\gs^{2K-1})$ with
$O(\gs^{2K})$ remainder (the first unmatched order is the absent level $K+1$'s first
entry $\gs^{2(K+1)-2}=\gs^{2K}$; at $K=4$ this is the machine-checked slope-$8$
residual).
Moreover: the diagonals $d_n=T_n(\Psi^n)$ are \emph{forced} (linear response
$\mathrm{diag}/n!$, by (ii)); the recursion is triangular (order $n$ never disturbs
orders $<n$); and everything off-diagonal is presentation freedom that feeds forward
into the higher $A_m$ and is re-absorbed by the recursion --- the canonical
(block-composite, as at the sextic and septic orders) and fully-polarized towers both
close, with the same forced $d_6$ but different tensors and hence different forced
$d_7$. \emph{Proof:} (i) + (ii) + induction on $n$. All geometric vertices enter the cubic
action only through the forced Veronese diagonals; the running rule (\cref{ped:eq:runningV}) is now a theorem.

\subparagraph{(iv) Gauge at all orders.} Block-Cholesky off-shell: with the conditional
deviations $Y_k:=\Sig_k-\hat H_k^{-1}L_k\Sig_{k-1}$ ($Y_2=\Sig_2-\hat
H_2^{-1}\tilde u_2$),
\begin{equation}
	S = S(\text{saddles})+\sum_{k\ge2}(-a_k)\,\hat H_k(Y_k,Y_k)
	\quad\text{exactly},\qquad
	\delta Y_k = -\tfrac12\,\hat H_k^{-1}(\delta\hat H_k)\,Y_k,
	\label{ped:eq:dYall}
\end{equation}
and with these rotations $\delta S=\delta[S(\text{saddles})]$ holds exactly off-shell
at every truncation (machine: $9\times10^{-16}$; dropping any one rotation leaves
exactly $-a_kY_k^\top(\delta\hat H_k)Y_k$). Homogeneous, no independent parameter at
any level: the pattern found at the quartic to septic orders is a theorem for the whole
tower --- the auxiliary sector is saddle-tracking plus contractible kernel rotations, and
one expects it to be in the BV class of the homogeneous auxiliary-rotation law
(\cref{ch:gauge-bv}) at every $k$, which is not proved.

\subparagraph{(v) The $O(\gs^6)$ instance, settled en passant.} The machine run executes
\eqref{ped:eq:diagrec} for $n=6\dots9$: the $\Sig_4$ pair $\{(4,4),(4,3,1)\}$ closes
orders $\gs^6$ \emph{and} $\gs^7$ with forced diagonals (residual log-slope $8.02$);
the source $\tilde u_4(\Sig_3)$ is forced (zeroing it leaves exactly the $A_6$
deficit) and \emph{uses} $\Sig_3$
($\partial\Sig_4^*/\partial\Sig_3=\hat H_4^{-1}L_4\neq0$). The tower rule
$\Sig_k\leftrightarrow(\tilde V_{2k},\tilde V_{2k+1})$ holds at $k=4$ and, by the
theorem, at every $k$. What the theorem does \emph{not} cover: the free-partition
conjecture for $(3,3,2)$-type couplings at $n\ge8$ (they break the Jacobi form ---
non-Gaussian shape; expected free by the pattern of the septic order, unchecked), the
convergence in $\gs$ of the infinite continued fraction (the statement is order-by-order),
and the CFT analytics of the nested $\hat H_k$ chain (now one condition per level).


